
%% 
%% Copyright 2019-2024 Elsevier Ltd
%% 
%% This file is part of the 'CAS Bundle'.
%% --------------------------------------
%% 
%% It may be distributed under the conditions of the LaTeX Project Public
%% License, either version 1.3c of this license or (at your option) any
%% later version.  The latest version of this license is in
%%    http://www.latex-project.org/lppl.txt
%% and version 1.3c or later is part of all distributions of LaTeX
%% version 1999/12/01 or later.
%% 
%% The list of all files belonging to the 'CAS Bundle' is
%% given in the file `manifest.txt'.
%% 
%% Template article for cas-sc documentclass for 
%% double column output.

\documentclass[a4paper,fleqn,doubleblind]{cas-sc}

% If the frontmatter runs over more than one page
% use the longmktitle option.

%\documentclass[a4paper,fleqn,longmktitle]{cas-sc}

\usepackage[numbers]{natbib}
%\usepackage[authoryear]{natbib}
% \usepackage[authoryear,longnamesfirst]{natbib}

%%%Author macros
\def\tsc#1{\csdef{#1}{\textsc{\lowercase{#1}}\xspace}}
\tsc{WGM}
\tsc{QE}
%%%

% Uncomment and use as if needed
%\newtheorem{theorem}{Theorem}
%\newtheorem{lemma}[theorem]{Lemma}
%\newdefinition{rmk}{Remark}
%\newproof{pf}{Proof}
%\newproof{pot}{Proof of Theorem \ref{thm}}

% Imported Packages from main.tex
\usepackage{amsmath}
\usepackage{amssymb}
\usepackage{graphicx}
\usepackage{textcomp}
\usepackage[fontsize=11.5pt]{scrextend}
\usepackage{setspace}
\usepackage{algorithm}
\usepackage{tabularx}
\usepackage{array}
\usepackage{float}
\usepackage{enumitem}
\usepackage{algpseudocode}
\usepackage{amsfonts}
\usepackage{cuted}
\usepackage{listings}
\usepackage[utf8]{inputenc}
\usepackage[most, minted]{tcolorbox}
\usepackage{minted}
\usepackage{booktabs}
\usepackage{multirow}
\usepackage[table]{xcolor}
\usepackage{caption}
\usepackage{cleveref}
\newcommand{\revise}[1]{{\color{blue}#1}}
\newcommand{\RightComment}[1]{%
  \hfill\makebox[0pt][r]{\textcolor{gray}{\footnotesize // #1}}%
}

\usepackage{hyperref}
\usepackage{xurl}
\usepackage{xr-hyper}
% \externaldocument{supplementary}
\hypersetup{
    colorlinks,
    linkcolor={red!50!black},
    citecolor={blue!50!black},
    urlcolor={black!80!black}
}
% \usepackage{natbib}
%  \bibpunct[, ]{(}{)}{,}{a}{}{,}%
%  \def\bibfont{\small}%
%  \def\bibsep{\smallskipamount}%
%  \def\bibhang{24pt}%
%  \def\newblock{\ }%
%  \def\BIBand{and}%

\newcommand{\rev}[1]{\textcolor{red}{#1}}

\graphicspath{{./}{figs/}}
\geometry{margin=1in}

\captionsetup[figure]{skip=2pt}

\newcommand{\ApplyManuscriptSpacing}{%
  \setstretch{1.6}%
  \setdisplayskipstretch{0.5}
  \selectfont
}

\definecolor{keywordcolor}{rgb}{0.13, 0.45, 0.75}
\definecolor{commentcolor}{rgb}{0.5, 0.5, 0.5}
\definecolor{phasecolor}{rgb}{0.2, 0.5, 0.2} 

\newcommand{\keyword}[1]{\textcolor{keywordcolor}{\textbf{#1}}}
\newcommand{\phase}[1]{\Statex \textcolor{phasecolor}{\textsc{#1}}}
\algnewcommand{\SideComment}[1]{\hfill \textcolor{commentcolor}{\(\triangleright\) \textit{#1}}}


\begin{document}
\let\WriteBookmarks\relax
\def\floatpagepagefraction{1}
\def\textpagefraction{.001}
\ApplyManuscriptSpacing

% Short title
\shorttitle{MarketToM}    

% Short author
\shortauthors{Anonymous Author(s)}  

% Main title of the paper
\title [mode = title]{MarketToM: A Theory of Mind Framework for Modeling Latent Mental States in Stock Trend Prediction \vspace{-30pt}}  

% Title footnote mark
% eg: \tnotemark[1]
% \tnotemark[1] 

% Title footnote 1.
% eg: \tnotetext[1]{Title footnote text}
% \tnotetext[1]{} 

% \author[1]{}

% Address/affiliation
% \affiliation[1]{}

% Corresponding author text
% \cortext[1]{Corresponding author}

% Footnote text
% \fntext[1]{}

% For a title note without a number/mark
%\nonumnote{}

% Here goes the abstract
\begin{abstract} 
Effective stock trend prediction requires financial decision support systems that can transform market quotes and textual information into behaviorally meaningful decision inputs. Yet most existing approaches treat these data primarily as surface-level sentiment cues and do not explicitly infer the latent mental states that underlie market actions. This paper develops MarketToM, a Theory-of-Mind-based framework for stock trend prediction that operationalizes latent mental-state inference within a heterogeneous multi-agent architecture. Specifically, MarketToM models the market through Retail, Institutional, and Arbitrageur agents, each representing a distinct cognitive pathway, and leverages large language models to infer beliefs, intentions, and emotions from multimodal market information. The framework further supports strategy refinement through error-triggered inter-agent comparison, enabling later inference to reuse corrected higher-order reasoning patterns rather than relying on static one-shot inference. Experiments on three real-world financial benchmarks show that MarketToM generally outperforms statistical, deep learning, multimodal, and LLM-based baselines in both accuracy and Matthews Correlation Coefficient. In addition, case-based and ablation analyses suggest that first-order Theory of Mind helps capture key components of anomaly formation, whereas second-order reasoning adds value in modeling recursive belief dynamics. These findings indicate that explicitly modeling latent mental states can move financial DSS beyond surface sentiment extraction toward more behaviorally grounded forecasting and more interpretable support for decision-making in complex market environments.
\vspace{-10pt}
\end{abstract}

% Use if graphical abstract is present
%\begin{graphicalabstract}
%\includegraphics{}
%\end{graphicalabstract}

% Research highlights
% \begin{highlights}
% \item Operationalizes Theory of Mind into a structured cognitive framework for stock prediction through an explicit belief–intention–emotion–action pathway.
% \item Introduces a heterogeneous multi-agent framework with inter-agent learning, enabling adaptive second-order ToM rather than static one-shot reasoning.
% \item Unifies predictive accuracy, adaptive refinement, and mechanism-oriented explanation in a single financial forecasting artifact.
% \end{highlights}


% Keywords
% Each keyword is seperated by \sep
\begin{keywords}
Theory of Mind \sep Stock Trend Prediction \sep LLMs \sep Explainable AI \sep Financial Decision Support Systems
\end{keywords}

\maketitle
\ApplyManuscriptSpacing

% Main text
\section{Introduction}\label{sec:introduction}
Building effective decision support systems (DSS) for stock trend prediction remains difficult \citep{hagenau2013automated,du2024explainable, liu2024multimodal}. A central reason is that classical financial theories, especially the Efficient Market Hypothesis, assume rational decision makers (\textit{homo economicus}), whereas real market behavior often departs from this idealization \citep{fama1970efficient, shiller2003efficient}. As a result, even professional investors and advanced predictive systems struggle to generate excess returns consistently; recent SPIVA scorecards show that most active managers underperform their benchmarks across major markets \citep{fama2010luck, spivaglobal2024}. 

Part of this challenge lies in how to model market data: they treat prices, volumes, and news as outputs of a relatively stable process, limiting support under nonlinear, non-stationary, and volatile market conditions \citep{nam2019financial}. Recent studies already show that decision-relevant market signals extend beyond isolated document sentiment to multi-source sentiment dynamics, inter-firm causal influence, collective incidental emotion, and investor co-attention in online communities \citep{yang2024uncovering,zhan2025stock,shao2025revisiting}. The pattern ``Buy the Rumor, Sell the News'' illustrates the problem: a sentiment-driven DSS may read strong expected earnings as bullish even when the expectation is already priced in, so the eventual release triggers profit-taking and a price decline despite positive news \citep{peterson2002rumor, schmidt2020rumor}. Addressing these limitations requires a conceptual shift from surface-level sentiment extraction to cognitive modeling that treats market data as behavioral traces of market participants' mental states, as illustrated in \Cref{fig:paradigm_comparison}.

\begin{figure}[pos=t]
\centering
\includegraphics[width=0.9\linewidth]{figs/paradigm_comparison.pdf}
\caption{\linespread{0.85}\selectfont Conceptual evolution of stock trend prediction paradigms.}
\label{fig:paradigm_comparison}
\end{figure}

Theory of Mind (ToM), defined as the cognitive ability to infer latent mental states such as beliefs, intentions, and emotions from observable behavior, offers a promising way to address this limitation \citep{premack1978does, frith2005theory}. For DSS, ToM matters because it shifts modeling from surface signals to the cognitive processes generating them, offering a principled basis for representing bounded rationality, although applying computational ToM to high-dimensional financial markets remains difficult \citep{baker2017rational, ying2024grounding}.

This paper operationalizes ToM through large language models (LLMs), making cognitively informed DSS design more feasible. Prior work suggests that LLMs show emergent ToM-like capabilities, while multi-step reasoning and in-context learning make them promising components for DSS that must reason over heterogeneous signals \citep{kosinski2024evaluating, shao2025revisiting, xu2025enigmatom, wei2022chain, xie2022explanation}. Related studies further indicate that combining LLMs with structured reasoning frameworks can improve interpretability and uncertainty estimation \citep{feng2025bird, ying2023neuro}. 

Building on this foundation, we propose the \textit{Market Theory-of-Mind} (MarketToM) framework, a DSS architecture designed to address cognition-specific decision challenges. MarketToM models three cognitive skills central to trading performance---fluid intelligence, cognitive reflection, and ToM \citep{corgnet2018makes}---and integrates heterogeneous market data through a ToM-guided Causal Cognitive Network (CCN), an LLM-based reasoning backend, and a Cognitive Enhancement Plugin (CEP) with backward inference to refine its strategy memory after observed prediction failures. Across three financial datasets, MarketToM achieves the strongest overall ACC and MCC among competing methods.

\revise{
Our framework contributes to DSS research in two ways. First, it provides design knowledge for AI-enabled financial DSS by operationalizing ToM as a heterogeneous multi-agent architecture. The proposed CCN organizes belief, intention, emotion, and action into an explicit cognitive reasoning pathway, while the CEP and backward inference mechanism allow agents to refine strategies after prediction failures. Second, it extends stock-prediction DSS beyond surface sentiment modeling by representing latent beliefs, intentions, and emotions as intermediate decision variables. This framing connects predictive analytics with mechanism-oriented explanation and helps explain market patterns such as why ``good news'' may already be priced in and why strategic profit-taking can emerge \citep{AbbasiEtAl2024, Tetlock2007}.
}


% Numbered list
% \section{Literature Review}\label{sec:related_work}
\section{Literature Review}\label{sec:related_work}

This paper integrates three streams: AI-driven stock prediction, computational ToM, and LLM-enabled financial DSS. For a recent overview of the shift from traditional models to LLM-based approaches, see \citet{darwish2025stock}.

Classic stock-prediction research relies on statistical models, technical indicators, and fundamental analysis, often assuming that structure in historical prices can forecast future movements \citep{fama1970efficient,schumaker2012evaluating,hagenau2013automated}. Yet many such regularities weaken after discovery, reflecting nonstationarity and market adaptation \citep{mclean2016does}. This motivates artifacts that combine predictive accuracy with explanation \citep{shmueli2011predictive, fernandez2022explaining}. Accordingly, AI research has expanded from technical signals to financial sentiment analysis (FSA), structure-aware graph neural network, and LLM-based text analysis \citep{du2023incorporating, deng2023llms,du2025modeling}. Prior studies show that media tone, online discussion, and investor sentiment are informative for market outcomes \citep{Tetlock2007, antweiler2004}, while domain lexicons and learning-based methods improve textual measurement \citep{loughran2011liability, frankel2022disclosure}. More recent work models text through networks and narratives rather than isolated words \citep{hoberg2016text, bybee2023narrative}. Even so, the dominant DSS pipeline still maps surface text signals to price movements, with little focus on explicitly inferring the latent beliefs, intentions, or strategies behind market actions \citep{schumaker2012evaluating,kraus2017decision}.

\begin{figure}[pos=t]
\centering
\includegraphics[width=0.8\columnwidth]{figs/false_belief_task.pdf}
\caption{\linespread{0.85}\selectfont The False Belief Task, a standard paradigm for assessing ToM.}
\label{fig:false_belief_task}
\end{figure}

Theory of Mind (ToM) is the capacity to infer others' beliefs, intentions, and emotions from observable behavior \citep{premack1978does, frith2005theory}. As illustrated in \Cref{fig:false_belief_task}, the False Belief Task captures this idea by testing whether one recognizes that another agent may act on a belief inconsistent with reality. For DSS, ToM is relevant because it provides a cognitive lens for modeling how decision makers interpret signals, anticipate others, and act under bounded rationality.

In market settings, ToM helps investors extract information from prices and coordinate expectations \citep{bruguier2010exploring, bao2024reading}. It is also associated with better use of market-order information, although trading profits depend on combining ToM with cognitive reflection and analytical capability \citep{corgnet2018makes, hefti2025mental}. At the same time, mentalizing can amplify bubble riding and risk contagion \citep{de2013mind, bossaerts2019perception}. These findings establish ToM as central to financial decision making, but the literature largely measures mentalizing intensity rather than distinguishing first-order from second-order recursion.

As illustrated in \Cref{fig:tom_type}, computational modeling of ToM offers a formal basis for reasoning about latent mental states \citep{baker2017rational}. Prior models vary in the mental factors and causal pathways they represent \citep{saritacs2025systematic, baker2011bayesian}, but their direct application to high-dimensional financial markets is difficult. LLMs provide a tractable alternative: recent studies report ToM-like social reasoning capabilities, though debate remains over whether they exhibit genuine mentalizing or only ToM-like behavior \citep{kosinski2024evaluating, strachan2024testing}. Prompting and symbolic scaffolds can improve LLM ToM performance \citep{wilf2024think, sclar2023minding}, and recent reviews further note that prompt-based procedures are widely used to elicit or enhance LLM mental-state reasoning \citep{chen2025theory}. A holistic evaluation framework for situated ToM further argues that LLM ToM evaluation should cover mental-state dimensions such as belief, intention, and emotion/affect \citep{ma2023towards}. These advances, however, have mostly been demonstrated in controlled settings rather than open-ended financial DSS contexts.

\begin{figure}[pos=t]
\centering
\includegraphics[width=1.0\columnwidth]{figs/tom_type.pdf}
\caption{\linespread{0.85}\selectfont Representative causal structures in computational ToM.}
\label{fig:tom_type}
\end{figure}

A related line of work embeds cognitive theories into agent architectures. Heterogeneous interacting agents have long been central to finance \citep{delong1990noise}, and recent systems such as TwinMarket use LLM-driven BDI agents to reproduce stylized market facts in simulations \citep{yuzhe2025twinmarket}. In parallel, frameworks such as CausalStock combine multimodal signals with causal discovery to improve interpretability at the variable level \citep{li2024causalstock}. These approaches are important, but they emphasize simulation fidelity or variable-level explanation rather than cognitive interpretability. MarketToM instead supports DSS design by embedding the LLM in a structured cognitive architecture that models multi-step market reasoning, second-order mental recursion, and error-driven strategy refinement.

\section{MarketToM Framework}\label{sec:method}
Financial markets are complex adaptive systems in which aggregate trends emerge from interactions among heterogeneous actors and shared narratives \citep{kahneman2011thinking}. As depicted in \Cref{fig:framework}, MarketToM operationalizes this view as a computational process for inferring the market's collective mental state through three components that map to cognitive skills linked to trading performance \citep{corgnet2018makes}: the LLM reasoning module corresponds to ``fluid intelligence,'' the Cognitive Enhancement Plugin (CEP) for error-triggered strategy refinement corresponds to ``cognitive reflection,'' and the ToM-guided CCN corresponds to ``Theory of Mind'' by reconstructing intentions and mental states from behavioral signals. This abstraction is consistent with prior work on market-level expectations and affective states \citep{de2013mind,bruguier2010exploring,bossaerts2019perception,hefti2025mental}.

\begin{figure}[pos=t]
\centering
\includegraphics[width=1.0\linewidth]{figs/framework.pdf}
\caption{\linespread{0.85}\selectfont An overview of MarketToM framework. }
\label{fig:framework}
\end{figure}


\subsection{Dependency Structure of the CCN}
The CCN models the market's collective cognitive process as a directed acyclic graph grounded in computational ToM, which explains action through probabilistic inference over latent mental states \citep{baker2011bayesian,baker2017rational}. In this literature, the graph direction is intended to encode asymmetric dependencies, not interchangeable associations: perceptions or evidence shape beliefs, beliefs together with desires/goals support an efficient plan of action (intention), and the plan guides action \citep{ho2022planning}. MarketToM adapts this logic to financial prediction by treating observable quote and text evidence as shaping beliefs, beliefs as guiding action-oriented intentions, environmental evidence and beliefs as jointly shaping emotions, and intentions and emotions as informing the next-step action prediction. This yields the dependency structure shown in the ``Causal Cognitive Network'' subgraph of \Cref{fig:framework}: $\mathbf{ES} \to b$, $b \to i$, $\mathbf{ES}, b \to e$, and $i, e \to a$.

Unlike classical ToM models that separately represent goals or desires, MarketToM treats profit maximization as a shared structural incentive among market participants, not as a time-varying latent state to be inferred at each prediction step. The model therefore focuses on the more dynamic states of belief, intention, and emotion, which mediate short-term market actions. The emotion pathway follows appraisal theory \citep{smith1993appraisal, lazarus1991progress}, under which affective states arise from evaluating new events in light of existing beliefs. This pathway is important in financial markets because risk perception, fear, optimism, and sentiment-driven reactions can shape how investors translate information and beliefs into actions.

\begin{table*}[t]
\centering
\small
\renewcommand{\arraystretch}{1.08}
\caption{Operational components and dependencies of the Causal Cognitive Network in MarketToM.}
\label{tab:ccn_components}
\begin{tabularx}{\textwidth}{>{\raggedright\arraybackslash}p{2.8cm}>{\raggedright\arraybackslash}X}
\toprule
\textbf{Component} & \textbf{Operational role in MarketToM} \\
\midrule
\multicolumn{2}{l}{\textit{Node set $V$}} \\
$\mathbf{ES}$ & Observable environmental states, including price information and public textual signals. \\
$b$ (belief) & Inferred assessment of market conditions and their likely implications. \\
$i$ (intention) & Action-oriented trading tendency derived from belief. \\
$e$ (emotion) & Affective appraisal inferred from environmental evidence and belief. \\
$a$ (action) & Predicted next-step market action, operationalized as the stock-trend label. \\
\midrule
\multicolumn{2}{l}{\textit{Edge set $E$}} \\
$\mathbf{ES} \to b$ & Observable market evidence shapes the inferred belief. \\
$b \to i$ & Belief is translated into an action-oriented trading tendency. \\
$\mathbf{ES}, b \to e$ & Emotion is inferred as an appraisal of environmental evidence under the current belief. \\
$i, e \to a$ & Intention and emotion jointly inform the predicted action. \\
\bottomrule
\end{tabularx}
\end{table*}


\subsection{Problem Formulation}
The task is next-day stock trend prediction under a $w$-step look-back window. At time $t$, the observation is $\mathbf{ES}_{t}^{w} = (\mathbf{es}_{t-w+1}, \dots, \mathbf{es}_{t})$, where each environmental state $\mathbf{es}_{t}$ contains quantitative quote data ($\text{quote}_t$) and public textual data ($\text{text}_t$). The target $a_{t+1}$ is a binary market-action label, with $a_{t+1}=1$ (Buy) for a predicted price increase and $a_{t+1}=0$ (Sell) for a decrease. The latent market mental state is $\mathcal{M}_t=\{b_t,i_t,e_t\}$, denoting belief, intention, and emotion.

\revise{Table~\ref{tab:ccn_components} summarizes how these latent states are operationalized within the CCN. This structural definition separates MarketToM from sentiment or narrative modeling by specifying belief, intention, and emotion as intermediate states with explicit parent conditions and downstream roles in the prediction path.}

The predictive probability $P(a_{t+1} \mid \mathbf{ES}_{t}^{w})$ is obtained by aggregating over latent mental states:

\begin{equation}
    P(a_{t+1} \mid \mathbf{ES}_{t}^{w}) = \sum_{\mathcal{M}_t} P(a_{t+1} \mid \mathcal{M}_t)P(\mathcal{M}_t \mid \mathbf{ES}_{t}^{w}).
  \label{eq:main_modeling_goal}
\end{equation}

Because exact marginalization in Eq.~\ref{eq:main_modeling_goal} requires summing over a large latent-state space and no closed-form parameterization of $P(\mathcal{M}_t \mid \mathbf{ES}_{t}^{w})$ is available for mixed textual, time-series, and cognitive signals, MarketToM adopts a heuristic MAP-style path approximation along the CCN dependency path:
\begin{equation}
  b_t^* = \underset{b_t}{\arg\max} \, P(b_t \mid \mathbf{ES}_t^w), 
  i_t^* = \underset{i_t}{\arg\max} \, P(i_t \mid b_t^*), 
  e_t^* = \underset{e_t}{\arg\max} \, P(e_t \mid b_t^*, \mathbf{ES}_t^w).
\end{equation}

From the inferred intention $i_t^*$ and emotion $e_t^*$, the upward probability $P_{\text{up}} \approx P(a_{t+1}=1 \mid i_t^*, e_t^*)$ is computed to determine the final action $\hat{a}_{t+1}$: $\hat{a}_{t+1} = \mathbb{I} \left( P_{\text{up}} > 0.5 \right)$, where $\mathbb{I}(\cdot)$ is the indicator function.

MarketToM approximates this sequential path by prompting an LLM to traverse the CCN. Prior work has used LLMs for probabilistic estimation~\citep{liu2025dellma, feng2025bird, requeima2024llm}, and theoretical studies suggest that sufficiently pretrained LLMs can perform implicit causal reasoning when conditioned on structured context~\citep{xie2022explanation, prystawski2023think, zhang2025and}. In each state-specific step, the prompt supplies the parent mental-state descriptions $\psi(\mathrm{Pa}(s_t))$ defined by the CCN, and the LLM outputs the most plausible child state $s_t$. \revise{MarketToM uses XML-style structured prompts to separate the role definition, CCN theory, task instruction, input fields, reasoning steps, and output schema; this design is intended to make the inference interface explicit and reproducible, and its robustness is examined through the prompt-format sensitivity diagnostic in Section~\ref{sec:ablation_study}~\citep{puerto2024code, li2025structured, mishra2023prompting}.}


\subsection{Forward Inference: Predicting Stock Trends}
The forward inference process implements the theory as a computational procedure that follows the forward inference paradigm. It proceeds in two stages: a dynamic reasoning stage to infer latent mental states and a probabilistic prediction stage to determine the action.


\subsubsection{Dynamic Mental State Reasoning}
The CEP is a dynamic repository of financial reasoning strategies. In this context, a ``strategy'' refers to a context-specific reasoning pattern stored as a state-strategy pair $(\psi(\text{Pa}(s_t)), \pi_{s,j})$, where $\psi(\text{Pa}(s_t))$ is the parent-state context, such as descriptions of beliefs or environmental states, that conditions the application of $\pi_{s,j}$, a textual guideline for inferring the target latent state. Strategies are organized in state-specific collections $\mathcal{D}_s$ for each latent state $s$.

In MarketToM, strategies are further specialized by agent role. For each agent role $k$ (Retail, Institutional, Arbitrageur) and each latent mental state $s \in \{b,i,e\}$, the CEP maintains a set of candidate strategies $\mathcal{D}_{s,k}$, and multiple strategies can be retrieved for the same role--state pair. Crucially, each strategy is designed to be ToM-aware: it not only provides guidance for inferring the target state under the CCN parents, but also embeds role-specific mental models of the other two agent roles, including both first-order cues (what other agent roles are likely to believe/intend/feel) and second-order cues (what other agent roles are likely to believe about each other and about the current role).

For a given agent role $k$ and target latent state $s_t$, the parent descriptions $\psi(\text{Pa}(s_t))$ form a query to retrieve the single most similar strategy $\pi_{s,k}^*$ from the CEP. Cosine similarity is computed between the query embedding, derived from $\psi(\text{Pa}(s_t))$, and the stored strategy embeddings in $\mathcal{D}_{s,k}$, followed by selecting the highest-similarity item. The reasoning step is:
\begin{equation}
s_{t,k}^* = \text{LLM}_{\text{Infer}}(\psi(\text{Pa}(s_t)), \pi_{s,k}^*)
\label{eq:state_inference}.
\end{equation}
The parent set $\text{Pa}(s_t)$ defined per the CCN structure (e.g., $\text{Pa}(i_t) = \{b_t^*\}$, while $\text{Pa}(e_t) = \{b_t^*, \mathbf{ES}_t^w\}$). Consequently, environmental evidence is incorporated only where the CCN explicitly specifies a direct dependency between the environmental states and the target latent state.

\revise{At the implementation level, each state-specific LLM call is required to return a JSON object with the agent role, target state type, and state description, and the response is parsed through a Pydantic schema before being stored as the corresponding belief, intention, or emotion field. Malformed or schema-inconsistent outputs are rejected and retried, so only validated state descriptions are propagated to downstream inference steps. This schema validation constrains the output interface and state-slot assignment, while the semantic validity of the inferred mental states is further examined through the external consistency and dependency-ablation analyses.}


\subsubsection{Heterogeneous Multi-Agent Framework and Dynamic Aggregation}
Accurately capturing uncertainty is essential in financial forecasting, where noisy, incomplete, and rapidly shifting information makes single deterministic predictions unreliable. Prior work shows that LLMs’ verbal probabilities are often poorly calibrated \citep{ma2025estimating}, frequently exhibiting overconfidence \citep{xiong2024can, kadavath2022language} and failing to reflect true predictive uncertainty \citep{feng2025bird}. This limitation is especially problematic in finance: market narratives are fragmented, sentiment signals are weak or contradictory, and contextual cues in textual data are inherently ambiguous. Under such conditions, relying on free-form verbalized probabilities is insufficient. Therefore, each agent is required to directly output an action label from \{Buy, Sell\}, and the model’s uncertainty is read from the logits assigned to these candidate labels. In this terminology, action-label probabilities are the actionable output, logits are the pre-softmax evidence over candidate labels, and $C_k$ is the log-confidence of the selected action label used for dynamic weighting.

\begin{figure}[pos=t]
\centering
\includegraphics[width=0.9\columnwidth]{figs/logits.pdf}
\caption{\linespread{0.85}\selectfont Schematic illustration of how contextual strength influences LLM logits and resulting softmax probabilities.}
\label{fig:logits}
\end{figure}

As illustrated in \Cref{fig:logits}, logits provide a more direct view of the model's internal evidence before softmax normalization. Strong and coherent contextual signals produce widely separated logits and sharply peaked action probabilities, whereas ambiguous or sentimentally mixed contexts produce logits of similar magnitude and flatter probability distributions. Because financial texts are often short, incomplete, or contradictory, this logit-level uncertainty is useful for distinguishing confident directional signals from weak or noisy ones.

MarketToM uses this uncertainty signal within a \textit{Heterogeneous Multi-Agent Framework} rather than relying on a single generic predictor. The market is modeled through three agent roles that mirror real participants: \textit{Retail agents}, who react intuitively to sentiment and trends; \textit{Institutional agents}, who emphasize fundamental value; and \textit{Arbitrageurs}, who look for mispricing opportunities.

Predictions from different agents are combined using dynamic weighted voting with two factors: (1) \textit{EMA Track Record} ($A_{k,t}$), and (2) \textit{Action-Label Log-Confidence} ($C_k$), where $C_k$ is computed from the direct Buy/Sell label probabilities obtained from logits via softmax, and the EMA memory length is controlled by $\beta \in [0,1)$. Let $\mathbf{A}_t = (A_{1,t}, \dots, A_{|\mathcal{K}|,t})^\top$ and $\mathbf{C} = (C_1, \dots, C_{|\mathcal{K}|})^\top$.

The aggregation weight vector is computed as:
\begin{equation}
\mathbf{W} = \text{Softmax} \left( \frac{\alpha \mathbf{A}_t + \gamma \mathbf{C}}{\tau} \right),
\label{eq:dynamic_weighting}
\end{equation}
where
$p_{\text{up},k} = P(a_{t+1}=\text{Buy}\mid \text{ctx}_k), \qquad
\hat{a}_{k,t+1} = \arg\max_{a \in \{\text{Buy}, \text{Sell}\}} P(a\mid \text{ctx}_k)$,
and $C_k = \log P(\hat{a}_{k,t+1}\mid \text{ctx}_k)$. Here, $p_{\text{up},k}$ is the agent-specific upward probability implied by the direct label distribution, $C_k$ is the log-confidence of the selected action label, and $\tau > 0$ is the softmax temperature used only for dynamic aggregation across agents. The scalar weight for agent role $k$ is the $k$-th element of $\mathbf{W}$, denoted by $W_k$. Lower $C_k$ leads to lower weight. The final upward probability is $P_{\text{up}} = \sum_k W_k \cdot p_{\text{up}, k}$, and the action $\hat{a}_{t+1}$ is determined by the 0.5 threshold. The full procedure is given in Algorithm~\ref{alg:parallel_dynamic_aggregation}.

\begin{algorithm}[t]
\caption{Agent Modeling and Dynamic Aggregation}
\label{alg:parallel_dynamic_aggregation}
\small
\begin{algorithmic}[1]
    \Require Environmental state $\mathbf{ES}_t^w$; Agent roles $\mathcal{K}=\{Ret, Inst, Arb\}$.
    \Ensure Upward probability $P_{\text{up}}$.
    \phase{Phase 1: Agent Modeling}
    \State \keyword{for} $k \in \mathcal{K}$ \keyword{do}
        \State \quad Retrieve Strategy $S_k$ from CEP
        \State \quad Construct context $\text{ctx}_k$ using $S_k$ and $\mathbf{ES}_t^w$
    \State \quad $p_{\text{up},k}, \hat{a}_{k,t+1}, C_k \leftarrow \text{LLM}(\text{ctx}_k; \{\text{Buy}, \text{Sell}\})$
    \State \keyword{end for}
    \phase{Phase 2: Dynamic Aggregation}
    \State Update EMA Track Record $A_{k,t}$
    \State $\mathbf{W} \leftarrow \text{Softmax}((\alpha \mathbf{A} + \gamma \mathbf{C}) / \tau)$
  \State $P_{\text{up}} \leftarrow \mathbf{W}^\top \mathbf{p}_{\text{up}}$
    \State \keyword{return} $P_{\text{up}}$
\end{algorithmic}
\end{algorithm}


\subsection{Backward Inference: Error-Triggered Strategy Refinement}
High computational capacity alone is insufficient for trading success; without cognitive reflection, traders are highly susceptible to behavioral biases such as conservatism, failing to effectively update their beliefs based on market signals \citep{corgnet2018makes}. Drawing on reflection-based feedback in LLM agent research, we operationalize this cognitive reflection through an \textit{Inter-Agent Strategy Refinement} mechanism tailored to the ToM setting. This module is not parametric learning with an explicit optimization objective or convergence guarantee. Instead, it updates the textual CEP strategy memory after feedback becomes available, so that later inference can reuse corrected reasoning patterns.

When a prediction error occurs, the realized label is used to partition the agents from the same forward pass into successful and failing sets:
\begin{equation}
\mathcal{S}_t=\{k\in\mathcal{K}: \hat{a}_{k,t+1}=a_{t+1}^{\text{actual}}\}, \qquad
\mathcal{F}_t=\{k\in\mathcal{K}: \hat{a}_{k,t+1}\neq a_{t+1}^{\text{actual}}\}.
\end{equation}
Successful agents are therefore defined entirely within the same forward pass: they are the roles whose individual predictions match the realized action, not external examples or future observations. For each failing agent $A \in \mathcal{F}_t$, the refinement prompt may include successful peers' inferred beliefs, intentions, and emotions; if $\mathcal{S}_t$ is empty, no peer trace is supplied.

Through this process, Agent $A$ integrates other-agent models into ToM-aware CEP updates. The update compares: (i) the strategies retrieved for Agent $A$ during forward inference, and (ii) the available successful agents' mental-state outputs that yielded the correct action. Based on this comparison, the refinement procedure decides whether the CEP should \emph{modify} an existing retrieved strategy (when a close but incomplete rule exists) or \emph{create} a new strategy (when the required inter-agent-conditioned logic is missing). Because each stored strategy already contains first-order and second-order other-agent mental models, this backward inference refines both levels: it strengthens the failing agent's direct modeling of others (first-order ToM) and its recursive modeling of how others model each other and the failing agent (second-order ToM).

The update rule is formulated as:
\begin{equation}
\Pi_{A}^{\text{updates}} = \text{LLM}_{\text{Refine}}(\text{Context}, \Pi_{A}^{\text{retrieved}}, \{\mathcal{M}_t^{k}\}_{k \in \mathcal{S}_t}, \text{Actual Action}).
\label{eq:strategy_update}
\end{equation}

\revise{The backward update avoids leakage by using only the original five-day input window, forward-generated mental-state traces, retrieved CEP strategies, predicted actions, and the realized label after the prediction horizon has passed. For benchmark comparisons, CEP updates occur only during training; validation and test evaluation freeze the CEP and use realized labels only for metrics, not strategy refinement. In Algorithm~\ref{alg:markettom_workflow}, this protocol is represented by \texttt{UpdateAllowed(phase, mode)}, which returns false for validation and test phases under the benchmark mode. In online deployment, CEP may be updated sequentially only after outcomes become observable, and each update affects only subsequent predictions.}

Algorithm \ref{alg:markettom_workflow} summarizes how Forward Inference and Backward Inference interact in the full workflow.

\begin{algorithm}[htbp]
\caption{The MarketToM Framework: Modeling and Evolution Cycle}
\small
\label{alg:markettom_workflow}
\begin{algorithmic}[1]
\Require Environmental state sequence $\mathbf{ES}_t^w$, Strategy Library $CEP$, Agent Roles $\mathcal{K}$, $CCN$.
\Ensure Predicted market action $\hat{a}_{t+1}$.
\Function{\keyword{ForwardInference}}{$\mathbf{ES}_t^w, CEP, \mathcal{K}, CCN$}
  \State $SP_t \gets \emptyset$
    \phase{Phase 1: Parallel Agent Modeling}
    \For{each Agent $k \in \mathcal{K}$}
        \State $\mathcal{M}_{t,k}^* \gets \emptyset$
        \For{each latent state $s_t \in \{b_t, i_t, e_t\}$ in $CCN$}
            \State $\text{Pa}(s_t) \gets CCN$.getParents($s_t, \mathcal{M}_{t,k}^*, \mathbf{ES}_t^w$)
            \State $\pi_{s,k}^* \gets \keyword{RetrieveStrategy}(\psi(\text{Pa}(s_t)), CEP, k)$
            \State $s_{t,k}^* \gets \keyword{LLM\_Infer}(\psi(\text{Pa}(s_t)), \pi_{s,k}^*)$
            \State $\mathcal{M}_{t,k}^*$.add($s_{t,k}^*$)
        \EndFor
    \State $(p_{\text{up},k}, \hat{a}_{k,t+1}, C_k) \gets \keyword{ActionPrediction}(\mathcal{M}_{t,k}^*, \mathbf{ES}_t^w, k)$
    \State $SP_t$.add($k, \mathcal{M}_{t,k}^*, p_{\text{up},k}, \hat{a}_{k,t+1}, C_k$)
    \EndFor
    \phase{Phase 2: Generate Prediction}
    \State $P_{\text{up}} \gets \keyword{DynamicAggregation}(SP_t, \tau)$
    \State $\hat{a}_{t+1} \gets \mathbb{I}(P_{\text{up}} > 0.5)$
    \State \Return $\hat{a}_{t+1}, SP_t$
\EndFunction
\Function{\keyword{BackwardInference}}{$a_{t+1}^{\text{actual}}, SP_t, CEP$}
  \State $\mathcal{F}_t$, $\mathcal{S}_t \gets \keyword{IdentifySets}(SP_t, a_{t+1}^{\text{actual}})$
  \For{each Agent$_A \in \mathcal{F}_t$}
    \State $\Pi_{A}^{\text{retrieved}} \gets SP_t$.getRetrievedStrategies(Agent$_A$)
    \State $\Pi_{A}^{\text{updates}} \gets \keyword{LLM\_Refine}(\text{Context}, \Pi_{A}^{\text{retrieved}}, \{\mathcal{M}_t^{k}\}_{k \in \mathcal{S}_t}, a_{t+1}^{\text{actual}})$
    \State $CEP$.applyUpdates(Agent$_A$, $\Pi_{A}^{\text{updates}}$)
  \EndFor
\EndFunction
\Procedure{\keyword{RunMarketToM}}{$\mathbf{ES}_t^w, phase, mode, ...$}
    \State $\hat{a}_{t+1}, SP_t \gets \keyword{ForwardInference}(\mathbf{ES}_t^w, ...)$
    \If{$a_{t+1}^{\text{actual}}$ is known \textbf{and} $\hat{a}_{t+1} \neq a_{t+1}^{\text{actual}}$ \textbf{and} $\keyword{UpdateAllowed}(phase, mode)$}
        \State $\keyword{BackwardInference}(a_{t+1}^{\text{actual}}, SP_t, CEP)$
    \EndIf
    \State \Return $\hat{a}_{t+1}$
\EndProcedure
\end{algorithmic}
\end{algorithm}


\revise{
\subsection{Computational Complexity and LLM Cost}
\label{sec:complexity_cost}
Because MarketToM uses LLM calls as part of its reasoning process, its operational cost is better characterized by both the number of LLM invocations and the associated token consumption. \Cref{tab:complexity_cost} summarizes the main cost sources under the default setting with three agent roles, three latent mental states, and one action-prediction call per agent.

\begin{table}[htbp]
\centering
\small
\caption{Computational complexity and LLM cost of MarketToM per prediction instance.}
\label{tab:complexity_cost}
\begin{tabular}{lcc}
\toprule
\textbf{Component} & \textbf{LLM calls} & \textbf{Token-cost source} \\
\midrule
Dynamic mental-state reasoning & $3 \times 3 = 9$ & Market window, CEP strategy, state prompt \\
Action prediction & $3 \times 1 = 3$ & Intention, emotion, and action-prediction prompt \\
Forward Inference & $12$ & Sum of the two rows above \\
Backward Inference & $0$--$3$ & Error-triggered prompt when CEP updates are allowed \\
\midrule
Total per instance & $12$--$15$ & Input and output tokens of all calls \\
\bottomrule
\end{tabular}
\end{table}

As shown in \Cref{tab:complexity_cost}, the fixed Forward Inference cost is 12 LLM calls per prediction instance. Backward Inference is not mandatory for every instance; it is triggered only when CEP updates are allowed and the aggregated prediction is wrong, adding at most one additional LLM call for each failing agent. Therefore, if $N_{\mathrm{eval}}$ prediction instances are evaluated and $\mathcal{F}_t$ is the set of failing agents for instance $t$ in update-enabled phases, the total number of LLM calls is
\begin{equation}
C_{\mathrm{call}} = N_{\mathrm{eval}}K_{\mathrm{agent}}(L_{\mathrm{state}}+1) + \sum_{t=1}^{N_{\mathrm{eval}}} |\mathcal{F}_t|.
\end{equation}
In the default configuration, this becomes $12N_{\mathrm{eval}}+\sum_t|\mathcal{F}_t|$. The corresponding token cost is
\begin{equation}
C_{\mathrm{token}} = \sum_{i=1}^{C_{\mathrm{call}}} (T_i^{\mathrm{in}} + T_i^{\mathrm{out}}),
\end{equation}
where $T_i^{\mathrm{in}}$ and $T_i^{\mathrm{out}}$ are the input and output tokens of the $i$-th LLM call. Monetary cost can be obtained by applying provider-specific input- and output-token prices. Non-LLM operations, such as CEP retrieval, use local sentence embeddings and cosine similarity and therefore do not contribute to API token cost. In deployment, cost can be reduced by disabling Backward Inference, shortening retrieved strategy text, caching embeddings, or using an open-source backbone.
}


\section{Experiments and Analysis}\label{sec:experiments}
This paper conducts a comprehensive empirical evaluation of MarketToM, encompassing benchmark comparisons, robustness and temporal generalization assessments, interpretability analyses, and ablation studies of its core components. \revise{The implementation, prompt template files, and example configuration are released in the project repository.\footnote{\url{https://github.com/yyveggie/MarketToM}}}


\subsection{Experimental Setup}
\subsubsection{Datasets and Preprocessing}
MarketToM is evaluated on three widely-used public datasets: {ACL18} \citep{xu2018stock} and {CMIN-US} \citep{luo2023causality}, two datasets of U.S. stocks paired with public textual data from different time periods; and {CMIN-CN} \citep{luo2023causality}, a corresponding dataset for the Chinese market.

However, the original textual data in these datasets exhibit considerable sparsity, with a missing rate of up to 99\% for certain stocks. To maintain data usability and cross-model comparability, stocks with extreme text sparsity are filtered out and evaluation is conducted on repeated random five-stock subsets drawn from the information-rich pool within each dataset. The filtering criterion is based on textual coverage rather than predictive performance. Specifically, for CMIN-CN, 11 stocks have no missing text and we retain stocks with a missing-text ratio below 2\%, resulting in 27 eligible stocks. For CMIN-US, 6 stocks have no missing text and the same below-2\% criterion results in 25 eligible stocks. Because all ACL18 stocks exhibit higher text sparsity, we retain stocks with a missing-text ratio below 15\%, resulting in 10 eligible stocks.

For each experimental run, five stocks are randomly selected from the corresponding eligible pool, and the same selected subset is used across all compared models. For the remaining few missing text entries, the placeholder ``No relevant news on this day'' is inserted to ensure structural consistency. Following prior studies using these datasets, the core stock trend prediction task is binary classification using a 5-day look-back window \citep{xu2022self, luo2023causality, zhang2024incorporating}. \revise{This choice is further supported by \citet{zhang2024incorporating}, who evaluate 5-, 7-, and 10-day windows and report that the 5-day window achieves the best ACC and MCC.} The prediction target is defined as a binary label for the next day's price movement: 1 if the closing price at day $t+1$ is higher than at day $t$, and 0 otherwise.

A systematic data-masking scheme is additionally applied to replace entity-specific cues in the textual data, including company, product, and CEO names, dates, years, store locations, and context-related comparable entities, with generic placeholders. This serves as a robustness and anti-leakage control by reducing reliance on explicit identifiers while preserving the semantic structure of the text.

The evaluation benchmarks MarketToM against four model families: classical Machine Learning (ML), Deep Learning (DL) time-series models, Multimodal Models (MM) including CMIN \citep{luo2023causality}, ACM24 \citep{du2024explainable}, Melody-GCN \citep{liu2024multimodal}, and CausalStock \citep{li2024causalstock}, and LLM-based methods including LLM-only and LLMFactor \citep{wang2024llmfactor}. \revise{The comparison uses the same underlying five-day quote--text information, next-day labels, repeated five-stock draws, train/validation/test splits, and held-out test instances for all methods. Because the baseline architectures differ, each method receives the modality representation supported by its original design, such as normalized quote features, text embeddings, graph inputs derived from the same window, or prompt-form quote--text context. No baseline uses future observations, external retrieval, or dataset-specific hyperparameter settings.}


\subsubsection{Metrics and Implementation Details}
Performance is evaluated using five standard metrics: Accuracy (ACC), Precision, Recall, F1 Score, and the Matthews Correlation Coefficient (MCC). The MCC is particularly emphasized for its robustness in handling imbalanced datasets. It considers all entries in the confusion matrix and provides a balanced measure even when class distributions differ substantially. It is defined as:
\begin{equation}
\text{MCC} = \frac{TP \times TN - FP \times FN}{\sqrt{(TP + FP)(TP + FN)(TN + FP)(TN + FN)}}.
\end{equation}

\revise{MarketToM is implemented with GPT-4o, with Qwen2.5-72B used for backbone sensitivity under the same split and protocol. LLM-dependent baselines also use GPT-4o without external retrieval; ML/DL baselines use normalized quote features and aligned text embeddings. Hyperparameters are selected on validation data using common search spaces across datasets, deterministic models use grid search, and stochastic models are averaged over five independent runs. Additional benchmark details and representative hyperparameters are provided in Appendix~\ref{app:benchmark_implementation}.}

\begin{table}[htbp]
\centering
\small
\renewcommand{\arraystretch}{0.8}
\caption{Comparison of Results on ACL18 Datasets}
\resizebox{\textwidth}{!}{%
\begin{tabular}{lllllll}
\toprule
\multicolumn{2}{c}{\textbf{Model Type}} & \textbf{ACC} $\uparrow$ & \textbf{Precision} $\uparrow$ & \textbf{Recall} $\uparrow$ & \textbf{F1 Score} $\uparrow$ & \textbf{MCC} $\uparrow$ \\
\midrule
\multirow{3}{*}{\textbf{ML}} & Gaussian NB & 0.5186 & 0.4107 & 0.3918 & 0.3374 & -0.0065 \\
 & RF & 0.5031 & 0.5004 & 0.4824 & 0.4573 & 0.0147 \\
 & Adaboost & 0.5397 & 0.5376 & 0.5303 & 0.4604 & 0.1124 \\
\midrule
\multirow{3}{*}{\textbf{DL}} & RNN & 0.5182±0.0037 & 0.4539±0.1399 & 0.4277±0.1660 & 0.4334±0.1526 & -0.0018±0.0780 \\
 & LSTM & 0.5132±0.0442 & 0.5623±0.1298 & 0.5543±0.2710 & 0.4501±0.2056 & 0.0307±0.0753 \\
 & GRU & 0.5261±0.0508 & 0.5702±0.1818 & 0.5287±0.2860 & 0.4515±0.1990 & 0.0567±0.0919 \\
\midrule
\multirow{5}{*}{\textbf{MM}} & CMIN & 0.5373±0.0052 & \textbf{0.6415±0.0369} & 0.4935±0.0434 & 0.5573±0.0375 & 0.0862±0.0111 \\
 & Melody-GCN & 0.5290±0.0132 & 0.5652±0.0918 & \textbf{0.5891±0.0537} & 0.5708±0.0472 & 0.0531±0.0182 \\
 & ACM24 & 0.5248±0.0156 & 0.4848±0.0440 & 0.5243±0.0923 & 0.4964±0.0338 & 0.0546±0.0309 \\
 & CausalStock & 0.5542±0.0045 & 0.5721±0.0510 & 0.5305±0.0480 & 0.5504±0.0410 & 0.1105±0.0380 \\
\midrule
\multirow{4}{*}{\textbf{LLM}} & LLM-only & 0.5195±0.0213 & 0.5237±0.0720 & 0.5003±0.0587 & 0.5056±0.0356 & 0.0474±0.0265 \\
 & LLMFactor & 0.5378±0.0538 & 0.5589±0.1016 & 0.4982±0.0987 & 0.5243±0.0953 & 0.0661±0.1018 \\
 & MarketToM (Qwen2.5-72B) & 0.5809±0.0513 & 0.6070±0.1171 & 0.5592±0.1079 & 0.5808±0.1078 & 0.1511±0.1160 \\
 & MarketToM (GPT-4o) & \textbf{0.5890±0.0486} & 0.6185±0.1093 & 0.5681±0.1032 & \textbf{0.5913±0.1006} & \textbf{0.1645±0.1082} \\
\bottomrule
\end{tabular}
}
\label{tab:comparison_results_acl18}
\end{table}

\begin{table}[htbp]
\centering
\small
\renewcommand{\arraystretch}{0.8}
\caption{Comparison of Results on CMIN-US Datasets}
\resizebox{\textwidth}{!}{%
\begin{tabular}{lllllll}
\toprule
\multicolumn{2}{c}{\textbf{Model Type}} & \textbf{ACC} $\uparrow$ & \textbf{Precision} $\uparrow$ & \textbf{Recall} $\uparrow$ & \textbf{F1 Score} $\uparrow$ & \textbf{MCC} $\uparrow$ \\
\midrule
\multirow{3}{*}{\textbf{ML}} & Gaussian NB & 0.4763 & 0.3461 & 0.6542 & 0.4488 & -0.0089 \\
 & RF & 0.5064 & 0.4802 & 0.3382 & 0.3184 & 0.0927 \\
 & Adaboost & 0.5237 & 0.5877 & 0.6837 & 0.5330 & 0.0836 \\
\midrule
\multirow{3}{*}{\textbf{DL}} & RNN & 0.5251±0.0249 & 0.5584±0.0600 & 0.7399±0.1629 & 0.5492±0.0912 & 0.0486±0.0666 \\
 & LSTM & 0.5058±0.0489 & 0.5660±0.1545 & \textbf{0.7723±0.2872} & 0.5558±0.2002 & 0.0260±0.0392 \\
 & GRU & 0.5098±0.0284 & 0.5673±0.0878 & 0.7483±0.1645 & 0.5572±0.0811 & 0.0283±0.0531 \\
\midrule
\multirow{5}{*}{\textbf{MM}} & CMIN & 0.5282±0.0143 & 0.4061±0.0949 & 0.5006±0.0154 & 0.4433±0.0532 & 0.0440±0.0272 \\
 & Melody-GCN & 0.5411±0.0054 & 0.5605±0.0642 & 0.4630±0.0158 & 0.5049±0.0254 & 0.0887±0.0248 \\
 & ACM24 & 0.5353±0.0030 & \textbf{0.6483±0.0812} & 0.6355±0.1046 & \textbf{0.6285±0.0086} & 0.0443±0.0168 \\
 & CausalStock & 0.5585±0.0120 & 0.6012±0.0550 & 0.5501±0.0420 & 0.5745±0.0350 & 0.1240±0.0310 \\
\midrule
\multirow{4}{*}{\textbf{LLM}} & LLM-only & 0.5180±0.0645 & 0.4964±0.0931 & 0.5726±0.0934 & 0.5271±0.0782 & 0.0486±0.1298 \\
 & LLMFactor & 0.5473±0.0289 & 0.6263±0.0493 & 0.5558±0.0630 & 0.5838±0.0136 & 0.1024±0.0577 \\
 & MarketToM (Qwen2.5-72B) & 0.5710±0.0142 & 0.6318±0.0847 & 0.5563±0.0543 & 0.5850±0.0321 & 0.1529±0.0362 \\
 & MarketToM (GPT-4o) & \textbf{0.5796±0.0131} & 0.6438±0.0793 & 0.5641±0.0508 & 0.5949±0.0298 & \textbf{0.1651±0.0339} \\
\bottomrule
\end{tabular}
}
\label{tab:comparison_results_cmin_us}
\end{table}

\begin{table}[htbp]
\centering
\small
\renewcommand{\arraystretch}{0.8}
\caption{Comparison of Results on CMIN-CN Datasets}
\resizebox{\textwidth}{!}{%
\begin{tabular}{lllllll}
\toprule
\multicolumn{2}{c}{\textbf{Model Type}} & \textbf{ACC} $\uparrow$ & \textbf{Precision} $\uparrow$ & \textbf{Recall} $\uparrow$ & \textbf{F1 Score} $\uparrow$ & \textbf{MCC} $\uparrow$ \\
\midrule
\multirow{3}{*}{\textbf{ML}} & Gaussian NB & 0.5364 & 0.5586 & \textbf{0.8020} & 0.6009 & 0.0691 \\
 & RF & 0.5091 & 0.4556 & 0.4688 & 0.3962 & 0.0384 \\
 & Adaboost & 0.5025 & 0.4748 & 0.6185 & 0.5056 & 0.0149 \\
\midrule
\multirow{3}{*}{\textbf{DL}} & RNN & 0.4880±0.0642 & 0.5831±0.2027 & 0.5410±0.4008 & 0.4171±0.2614 & 0.0095±0.0844 \\
 & LSTM & 0.4960±0.0571 & \textbf{0.6583±0.2141} & 0.5204±0.4010 & 0.3900±0.2612 & 0.0122±0.0415 \\
 & GRU & 0.5036±0.0666 & 0.5635±0.1480 & 0.7275±0.3643 & 0.5237±0.2200 & 0.0104±0.0457 \\
\midrule
\multirow{5}{*}{\textbf{MM}} & CMIN & 0.5218±0.0093 & 0.5593±0.0400 & 0.6159±0.0548 & 0.5826±0.0122 & 0.0327±0.0236 \\
 & Melody-GCN & 0.5301±0.0188 & 0.5401±0.0448 & 0.6146±0.0588 & 0.5705±0.0103 & 0.0634±0.0469 \\
 & ACM24 & 0.5191±0.0110 & 0.5848±0.0267 & 0.5386±0.0618 & 0.5577±0.0277 & 0.0332±0.0324 \\
 & CausalStock & 0.5492±0.0105 & 0.5710±0.0620 & 0.5920±0.0580 & 0.5813±0.0450 & 0.0915±0.0420 \\
\midrule
\multirow{4}{*}{\textbf{LLM}} & LLM-only & 0.5236±0.0818 & 0.5506±0.0998 & 0.5441±0.0554 & 0.5439±0.0706 & 0.0472±0.1643 \\
 & LLMFactor & 0.5333±0.0598 & 0.5423±0.0774 & 0.5492±0.1382 & 0.5428±0.1011 & 0.0607±0.1226 \\
 & MarketToM (Qwen2.5-72B) & 0.5647±0.0026 & 0.6144±0.0638 & 0.6629±0.0746 & 0.6304±0.0092 & 0.1203±0.0159 \\
 & MarketToM (GPT-4o) & \textbf{0.5732±0.0024} & 0.6261±0.0604 & 0.6735±0.0701 & \textbf{0.6417±0.0084} & \textbf{0.1342±0.0148} \\
\bottomrule
\end{tabular}
}
\label{tab:comparison_results_cmin_cn}
\end{table}


\subsection{Main Results and Analysis}
As shown in Tables~\ref{tab:comparison_results_acl18} to \ref{tab:comparison_results_cmin_cn}, all baselines are evaluated consistently. Across all three datasets, MarketToM attains the highest mean ACC and MCC, while F1 remains stable and competitive, though not uniformly the top value. Statistical significance is assessed with Pesaran--Timmermann directional accuracy tests and paired Diebold--Mariano tests, pooled over the held-out test predictions of each dataset. As reported in Table~\ref{tab:pesaran_timmermann_results}, both MarketToM variants reject the null hypothesis of independence between predicted and realized directions across all three datasets, with positive excess hit rates in every case. Table~\ref{tab:diebold_mariano_summary} further shows that MarketToM yields lower paired prediction losses than almost all compared baselines under both 0--1 and Brier loss, with many pairwise comparisons statistically significant, particularly on the larger CMIN-US and CMIN-CN test sets.

\begin{table}[htbp]
\centering
\small
\renewcommand{\arraystretch}{0.85}
\caption{Pesaran--Timmermann directional accuracy tests for MarketToM.}
\resizebox{\textwidth}{!}{%
\begin{tabular}{llccccc}
\toprule
\textbf{Dataset} & \textbf{Model} & \textbf{$N$} & \textbf{Hit Rate} & \textbf{Indep. Benchmark} & \textbf{Excess Hit} & \textbf{PT Statistic} \\
\midrule
ACL18 & MarketToM (GPT-4o) & 600 & 0.5812 & 0.5012 & 0.0800 & 3.9520$^{***}$ \\
ACL18 & MarketToM (Qwen2.5-72B) & 600 & 0.5748 & 0.5014 & 0.0734 & 3.6357$^{***}$ \\
CMIN-US & MarketToM (GPT-4o) & 1500 & 0.5805 & 0.5046 & 0.0759 & 6.0476$^{***}$ \\
CMIN-US & MarketToM (Qwen2.5-72B) & 1500 & 0.5725 & 0.5053 & 0.0672 & 5.4049$^{***}$ \\
CMIN-CN & MarketToM (GPT-4o) & 1512 & 0.5819 & 0.5175 & 0.0644 & 5.3793$^{***}$ \\
CMIN-CN & MarketToM (Qwen2.5-72B) & 1512 & 0.5691 & 0.5177 & 0.0514 & 4.2920$^{***}$ \\
\bottomrule
\end{tabular}
}
\begin{flushleft}
\footnotesize Note: The null hypothesis is that predicted and realized directions are independent. $N$ is the number of pooled held-out test predictions per dataset, computed over the full eligible stock pool. $^{*}$, $^{**}$, and $^{***}$ indicate two-sided $p<0.05$, $p<0.01$, and $p<0.001$, respectively.
\end{flushleft}
\label{tab:pesaran_timmermann_results}
\end{table}

\begin{table}[htbp]
\centering
\small
\renewcommand{\arraystretch}{0.85}
\caption{Summary of paired Diebold--Mariano tests for MarketToM against baselines.}
\resizebox{\textwidth}{!}{%
\begin{tabular}{lllcccc}
\toprule
\textbf{Dataset} & \textbf{Reference Model} & \textbf{Loss} & \textbf{Comparisons} & \textbf{Lower Loss} & \textbf{Sig. Lower Loss} & \textbf{Median Loss Diff.} \\
\midrule
ACL18 & MarketToM (GPT-4o) & 0--1 & 13 & 13 & 6 & 0.0612 \\
ACL18 & MarketToM (GPT-4o) & Brier & 13 & 13 & 6 & 0.0322 \\
ACL18 & MarketToM (Qwen2.5-72B) & 0--1 & 13 & 12 & 5 & 0.0548 \\
ACL18 & MarketToM (Qwen2.5-72B) & Brier & 13 & 12 & 4 & 0.0193 \\
CMIN-US & MarketToM (GPT-4o) & 0--1 & 13 & 13 & 8 & 0.0431 \\
CMIN-US & MarketToM (GPT-4o) & Brier & 13 & 13 & 8 & 0.0202 \\
CMIN-US & MarketToM (Qwen2.5-72B) & 0--1 & 13 & 12 & 7 & 0.0349 \\
CMIN-US & MarketToM (Qwen2.5-72B) & Brier & 13 & 12 & 6 & 0.0153 \\
CMIN-CN & MarketToM (GPT-4o) & 0--1 & 13 & 13 & 11 & 0.0585 \\
CMIN-CN & MarketToM (GPT-4o) & Brier & 13 & 13 & 10 & 0.0231 \\
CMIN-CN & MarketToM (Qwen2.5-72B) & 0--1 & 13 & 12 & 9 & 0.0455 \\
CMIN-CN & MarketToM (Qwen2.5-72B) & Brier & 13 & 12 & 7 & 0.0178 \\
\bottomrule
\end{tabular}
}
\begin{flushleft}
\footnotesize Note: Each row summarizes 13 paired Diebold--Mariano comparisons between the reference MarketToM model and the other models on the same dataset. The loss difference is baseline loss minus MarketToM loss, so a positive value indicates lower MarketToM loss. ``Sig. Lower Loss'' counts two-sided tests with $p<0.05$. These paired tests are used as complementary sample-level evidence for loss differences, rather than as a substitute for the broader robustness, temporal-generalization, and ablation analyses. Additional pairwise results against representative strong baselines are reported in Appendix~\ref{app:dm_tests}.
\end{flushleft}
\label{tab:diebold_mariano_summary}
\end{table}

\textit{Strong and Balanced Performance.} MarketToM achieves the highest mean ACC and MCC across all three datasets. Compared to the strongest external baseline (CausalStock), relative ACC gains are approximately $+6.28\%$, $+3.78\%$, and $+4.37\%$ on ACL18, CMIN-US, and CMIN-CN, respectively, while relative MCC gains are approximately $+48.87\%$, $+33.15\%$, and $+46.67\%$. These results highlight MarketToM's robustness, particularly under class imbalance where MCC serves as a more stringent indicator. F1 is consistently competitive across datasets, though it is not always the top value.

\textit{Architectural Advantage over Multimodal and LLM-based Methods.} The LLM-only baseline's performance frequently falls short of even non-LLM methods, suggesting that direct application of an LLM without structural reasoning is often insufficient in this setting. LLMFactor, which uses an LLM to automatically extract predictive text features, achieves notable gains. CausalStock, treated here as a multimodal-causal baseline, also improves performance by denoising news and modeling inter-stock statistical dependencies. However, MarketToM delivers additional improvements across datasets, indicating that multi-step cognitive reasoning provides a systematic architectural benefit beyond factor extraction or statistical graph learning alone.

\textit{Backbone Robustness (GPT-4o vs.\ Qwen2.5-72B).} Replacing GPT-4o with the open-source backbone Qwen2.5-72B in the same MarketToM pipeline still yields stronger ACC and MCC than non-MarketToM baselines across all three datasets. Relative to the strongest external baseline (CausalStock), the Qwen2.5-72B variant retains positive gains in ACC ($+4.82\%$, $+2.24\%$, $+2.82\%$) and MCC ($+36.74\%$, $+23.31\%$, $+31.48\%$) on ACL18, CMIN-US, and CMIN-CN, respectively. Meanwhile, the GPT-4o variant remains consistently higher than Qwen2.5-72B (about $+0.84$ percentage points in ACC and about $+0.0132$ in MCC on average), indicating stable headroom from backbone capacity. This pattern (GPT-4o $>$ Qwen2.5-72B, while both $>$ non-MarketToM baselines) suggests that the observed gains are not tied to a single proprietary model and that the proposed cognitive architecture transfers to at least one fully open-source backbone.

\textit{Effectiveness in Volatile Markets.} The Chinese stock market, unlike the U.S. market, exhibits higher volatility and stronger retail-driven sentiment dynamics. MarketToM’s MCC gain on the CMIN-CN dataset is consistent with the view that ToM-based reasoning may better capture collective sentiment and herd behavior in such environments.


\subsection{Robustness and Temporal Generalization}
We evaluate robustness through three targeted checks: temporal holdout performance on a later period, entity masking to reduce reliance on explicit identifiers, and semantic perturbations that preserve factual content while altering linguistic framing.

\textit{Temporal Generalization on an Unseen Time Period}
\begin{table}[htbp]
\centering
\small
\renewcommand{\arraystretch}{0.8}
\caption{Performance Comparison on a Temporal Holdout Split}
\resizebox{\textwidth}{!}{%
\begin{tabular}{l lllll}
\toprule
\textbf{Model} & \textbf{ACC} $\uparrow$ & \textbf{Precision} $\uparrow$ & \textbf{Recall} $\uparrow$ & \textbf{F1 Score} $\uparrow$ & \textbf{MCC} $\uparrow$ \\
\midrule
CMIN & 0.5355±0.0297 & 0.4982±0.0299 & 0.5282±0.0371 & 0.5127±0.0329 & 0.0698±0.0599 \\
Melody-GCN & 0.5409±0.0765 & 0.5036±0.0747 & 0.5471±0.0909 & 0.5241±0.0807 & 0.0827±0.1541 \\
ACM24 & 0.5326±0.0391 & 0.4944±0.0385 & 0.5429±0.0652 & 0.5172±0.0501 & 0.0665±0.0814 \\
CausalStock & 0.5589±0.0351 & 0.5211±0.0419 & 0.5481±0.0381 & 0.5341±0.0391 & 0.1181±0.0749 \\
LLMFactor & 0.5485±0.0412 & 0.5121±0.0529 & 0.5351±0.0611 & 0.5231±0.0479 & 0.0951±0.0819 \\
MarketToM (Qwen2.5-72B) & 0.5687±0.0556 & 0.5344±0.0594 & 0.5512±0.0476 & 0.5425±0.0527 & 0.1351±0.1098 \\
MarketToM (GPT-4o) & \textbf{0.5768±0.0523} & \textbf{0.5431±0.0562} & \textbf{0.5594±0.0449} & \textbf{0.5509±0.0496} & \textbf{0.1482±0.1027} \\
\bottomrule
\end{tabular}
}
\label{tab:temporal_holdout_split}
\end{table}

On the January--April 2025 temporal holdout, MarketToM (GPT-4o) achieves the best ACC ($0.5768$) and MCC in Table~\ref{tab:temporal_holdout_split}, exceeding the strongest non-MarketToM baseline, CausalStock, by $+0.0179$ ACC and $+0.0301$ MCC. This holdout period is also later than GPT-4o's October 2023 training cutoff reported by OpenAI.

\textit{Robustness to Entity Masking}
Entity masking is evaluated on 1,000 original/anonymized sample pairs. As shown in \Cref{tab:masked_unmasked}, predictions remain highly consistent (vector similarity $95.94\%$; Wasserstein distance $0.0406\pm0.0193$), while the small performance drop suggests that entity-specific priors provide only a weak contribution.

\textit{Robustness to Semantic Perturbations}
Semantic robustness is tested on ACL18 by introducing subtle LLM-generated framing changes while preserving factual and quantitative content (e.g., ``cautious optimism'' to neutral earnings reports); $\Delta$ in \Cref{tab:stress_test_results} denotes relative change.

\begin{table*}[!ht]
\centering
\small
\renewcommand{\arraystretch}{0.8}
\caption{Consistency of MarketToM Predictions With and Without Entity Masking}
\resizebox{\textwidth}{!}{%
\begin{tabular}{llllll}
\toprule
\textbf{Version} & \textbf{ACC} $\uparrow$ & \textbf{F1 Score} $\uparrow$ & \textbf{MCC} $\uparrow$ & \textbf{Vec. Similarity} $\uparrow$ & \textbf{Wass. Distance} $\downarrow$ \\
\midrule
Masked & 0.5616 ± 0.0173 & 0.5343 ± 0.0126 & 0.1240 ± 0.0138 & \multirow{2}{*}{0.9594 ± 0.0193} & \multirow{2}{*}{0.0406 ± 0.0193} \\
Unmasked & 0.5650 ± 0.0130 & 0.5571 ± 0.0089 & 0.1304 ± 0.0066 & & \\
\bottomrule
\end{tabular}
}
\label{tab:masked_unmasked}
\end{table*}

\begin{table*}[!ht]
\centering
\caption{\linespread{0.85}\selectfont Stress test results on the ACL18 dataset with semantic perturbations.}
\footnotesize
\setlength{\tabcolsep}{4pt}
\resizebox{\textwidth}{!}{%
\begin{tabular}{llccccc}
\toprule
\multicolumn{2}{c}{\textbf{Model Type}} & \textbf{ACC ($\Delta$)} $\uparrow$ & \textbf{Precision ($\Delta$)} $\uparrow$ & \textbf{Recall ($\Delta$)} $\uparrow$ & \textbf{F1 Score ($\Delta$)} $\uparrow$ & \textbf{MCC ($\Delta$)} $\uparrow$ \\
\midrule
\multirow{3}{*}{\textbf{ML}} & Gaussian NB & \shortstack{0.4564 \\ {\scriptsize -12.03\%}} & \shortstack{0.4509 \\ {\scriptsize +9.81\%}} & \shortstack{0.4808 \\ {\scriptsize +22.68\%}} & \shortstack{0.4546 \\ {\scriptsize +34.73\%}} & \shortstack{-0.0918 \\ {\scriptsize n/a (near-zero baseline)}} \\
 & RF & \shortstack{0.4786 \\ {\scriptsize -4.88\%}} & \shortstack{0.4789 \\ {\scriptsize -4.32\%}} & \shortstack{0.5656 \\ {\scriptsize +17.28\%}} & \shortstack{0.5117 \\ {\scriptsize +11.92\%}} & \shortstack{-0.0434 \\ {\scriptsize n/a (near-zero baseline)}} \\
 & Adaboost & \shortstack{0.5016 \\ {\scriptsize -7.12\%}} & \shortstack{0.5051 \\ {\scriptsize -6.08\%}} & \shortstack{0.3352 \\ {\scriptsize -36.82\%}} & \shortstack{0.3647 \\ {\scriptsize -20.77\%}} & \shortstack{0.0043 \\ {\scriptsize -96.18\%}} \\
\midrule
\multirow{3}{*}{\textbf{DL}} & RNN & \shortstack{0.4896$\pm$0.0174 \\ {\scriptsize -4.42\%}} & \shortstack{0.4891$\pm$0.0063 \\ {\scriptsize +8.93\%}} & \shortstack{0.6244$\pm$0.2259 \\ {\scriptsize +18.58\%}} & \shortstack{0.5145$\pm$0.0831 \\ {\scriptsize +10.43\%}} & \shortstack{-0.0268$\pm$0.0509 \\ {\scriptsize n/a (near-zero baseline)}} \\
 & LSTM & \shortstack{0.5099$\pm$0.0011 \\ {\scriptsize -1.32\%}} & \shortstack{0.5628$\pm$0.1203 \\ {\scriptsize -7.58\%}} & \shortstack{0.5444$\pm$0.3834 \\ {\scriptsize -42.87\%}} & \shortstack{0.4547$\pm$0.2507 \\ {\scriptsize -22.13\%}} & \shortstack{0.0439$\pm$0.0311 \\ {\scriptsize -43.68\%}} \\
 & GRU & \shortstack{0.5128$\pm$0.0053 \\ {\scriptsize -2.48\%}} & \shortstack{0.5108$\pm$0.0072 \\ {\scriptsize -10.42\%}} & \shortstack{\textbf{0.6696}$\pm$\textbf{0.2209} \\ {\scriptsize \textbf{+26.73\%}}} & \shortstack{\textbf{0.5663}$\pm$\textbf{0.0826} \\ {\scriptsize \textbf{+25.38\%}}} & \shortstack{0.0345$\pm$0.0214 \\ {\scriptsize -39.17\%}} \\
\midrule
\multirow{4}{*}{\textbf{MM}} & CMIN & \shortstack{0.5136$\pm$0.0141 \\ {\scriptsize -4.38\%}} & \shortstack{0.5199$\pm$0.0099 \\ {\scriptsize -18.92\%}} & \shortstack{0.4476$\pm$0.2726 \\ {\scriptsize -9.28\%}} & \shortstack{0.4488$\pm$0.1324 \\ {\scriptsize -19.53\%}} & \shortstack{0.0317$\pm$0.0066 \\ {\scriptsize -63.17\%}} \\
 & Melody-GCN & \shortstack{0.5166$\pm$0.0232 \\ {\scriptsize -2.32\%}} & \shortstack{0.5413$\pm$0.0395 \\ {\scriptsize -4.18\%}} & \shortstack{0.4476$\pm$0.3638 \\ {\scriptsize -23.98\%}} & \shortstack{0.4187$\pm$0.2018 \\ {\scriptsize -26.73\%}} & \shortstack{0.0506$\pm$0.0203 \\ {\scriptsize -4.68\%}} \\
 & ACM24 & \shortstack{0.5062$\pm$0.0026 \\ {\scriptsize -3.52\%}} & \shortstack{0.5062$\pm$0.0026 \\ {\scriptsize +4.42\%}} & \shortstack{0.5321$\pm$0.2861 \\ {\scriptsize +1.48\%}} & \shortstack{0.4886$\pm$0.1376 \\ {\scriptsize -1.62\%}} & \shortstack{0.0164$\pm$0.0112 \\ {\scriptsize -69.97\%}} \\
\midrule
\multirow{5}{*}{\textbf{LLM}} & LLM-only & \shortstack{0.5042$\pm$0.0227 \\ {\scriptsize -3.02\%}} & \shortstack{0.5032$\pm$0.0022 \\ {\scriptsize -3.88\%}} & \shortstack{0.6079$\pm$0.2738 \\ {\scriptsize +21.53\%}} & \shortstack{0.5259$\pm$0.1299 \\ {\scriptsize +4.03\%}} & \shortstack{0.0109$\pm$0.0072 \\ {\scriptsize -76.98\%}} \\
 & LLMFactor & \shortstack{0.5354$\pm$0.0232 \\ {\scriptsize -0.48\%}} & \shortstack{0.5633$\pm$0.0493 \\ {\scriptsize +0.82\%}} & \shortstack{0.4921$\pm$0.3295 \\ {\scriptsize -1.18\%}} & \shortstack{0.4661$\pm$0.1806 \\ {\scriptsize -11.13\%}} & \shortstack{0.0913$\pm$0.0166 \\ {\scriptsize +38.12\%}} \\
 & MarketToM (Qwen2.5-72B) & \shortstack{0.5732$\pm$0.0175 \\ {\scriptsize -1.28\%}} & \shortstack{0.6249$\pm$0.0856 \\ {\scriptsize +3.02\%}} & \shortstack{0.5016$\pm$0.2629 \\ {\scriptsize -10.32\%}} & \shortstack{0.5108$\pm$0.1293 \\ {\scriptsize -12.08\%}} & \shortstack{0.1716$\pm$0.0274 \\ {\scriptsize +13.63\%}} \\
 & MarketToM (GPT-4o) & \shortstack{\textbf{0.5823}$\pm$\textbf{0.0164} \\ {\scriptsize \textbf{-1.12\%}}} & \shortstack{\textbf{0.6392}$\pm$\textbf{0.0801} \\ {\scriptsize \textbf{+3.32\%}}} & \shortstack{0.5119$\pm$0.2497 \\ {\scriptsize -9.88\%}} & \shortstack{0.5213$\pm$0.1198 \\ {\scriptsize -11.83\%}} & \shortstack{\textbf{0.1874}$\pm$\textbf{0.0246} \\ {\scriptsize \textbf{+13.87\%}}} \\
 & MarketToM-NoCEP & \shortstack{0.5232$\pm$0.0026 \\ {\scriptsize -2.92\%}} & \shortstack{0.5159$\pm$0.0053 \\ {\scriptsize +0.78\%}} & \shortstack{0.5084$\pm$0.2199 \\ {\scriptsize -5.88\%}} & \shortstack{0.5194$\pm$0.0698 \\ {\scriptsize -1.32\%}} & \shortstack{0.0707$\pm$0.0326 \\ {\scriptsize -8.78\%}} \\
\bottomrule
\end{tabular}
}
\label{tab:stress_test_results}
\end{table*}

Table~\ref{tab:stress_test_results} shows that most baselines are sensitive to perturbation, with degraded ACC/MCC or unstable near-zero MCC behavior. MarketToM instead shows only a small ACC drop and higher MCC, suggesting more balanced classification under semantic ambiguity.


\revise{
\subsection{Trading Performance Evaluation}
\label{sec:trading_simulation}
Beyond ACC and MCC, we conduct an exploratory economic sanity check using a simple \emph{long--cash} rule: Buy predictions invest for day $t+1$, while Sell predictions remain in cash. Position changes incur $10$ bps round-trip transaction costs; each dataset--model pair is evaluated over $5$ seeds, with each seed covering approximately one trading year of daily predictions, and per-stock equity curves are averaged. Table~\ref{tab:trading_simulation} reports cumulative and annualized return, annualized volatility, Sharpe ratio, absolute maximum drawdown, and turnover.

\begin{table}[htbp]
\centering
\small
\renewcommand{\arraystretch}{0.8}
\caption{Trading-simulation results with $10$ bps round-trip transaction costs. Cum.\ Ret.: cumulative return; Ann.\ Ret.: annualized return; Ann.\ Vol.: annualized volatility; SR: Sharpe ratio; MDD: absolute maximum drawdown; Turn.: average turnover. All values are averaged over 5 seeds, each covering approximately one trading year. Returns are annualized by compounding the realized per-period returns to a standard $252$-trading-day year, volatility is the per-period standard deviation scaled by $\sqrt{252}$, and SR is computed as Ann.\ Ret.\ divided by Ann.\ Vol.\ with the risk-free rate set to zero.}
\label{tab:trading_simulation}
\resizebox{\textwidth}{!}{%
\begin{tabular}{lllrrrrrr}
\toprule
\textbf{Dataset} & \multicolumn{2}{c}{\textbf{Model Type}} & \textbf{Cum.\ Ret.} $\uparrow$ & \textbf{Ann.\ Ret.} $\uparrow$ & \textbf{Ann.\ Vol.} $\downarrow$ & \textbf{SR} $\uparrow$ & \textbf{MDD} $\downarrow$ & \textbf{Turn.} \\
\midrule
\multirow{14}{*}{ACL18}
 & \multirow{3}{*}{\textbf{ML}}
    & Gaussian NB & -0.0212 & -0.0238 & 0.2241 & -0.1062 & 0.1287 & 0.4625 \\
 && RF           & -0.0061 & -0.0069 & 0.2271 & -0.0304 & 0.1156 & 0.4921 \\
 && Adaboost     & 0.1236 & 0.1365 & 0.2191 & 0.6230 & 0.1046 & 0.4835 \\
\cmidrule(l){2-9}
 & \multirow{3}{*}{\textbf{DL}}
    & RNN  & -0.0184 & -0.0207 & 0.2068 & -0.1001 & 0.1213 & 0.4949 \\
 && LSTM & 0.0958 & 0.1113 & \textbf{0.1967} & 0.5658 & 0.0949 & 0.5025 \\
 && GRU  & 0.1655 & 0.1919 & 0.2069 & 0.9275 & \textbf{0.0748} & 0.5119 \\
\cmidrule(l){2-9}
 & \multirow{4}{*}{\textbf{MM}}
    & CMIN        & 0.1658 & 0.1970 & 0.2311 & 0.8524 & 0.0969 & 0.4765 \\
 && Melody-GCN  & 0.2158 & 0.2339 & 0.2334 & 1.0021 & 0.1073 & 0.4435 \\
 && ACM24       & 0.1781 & 0.2006 & 0.2221 & 0.9032 & 0.1021 & 0.4575 \\
 && CausalStock & 0.2249 & 0.2421 & 0.2261 & 1.0708 & 0.0859 & 0.4519 \\
\cmidrule(l){2-9}
 & \multirow{4}{*}{\textbf{LLM}}
    & LLM-only                 & 0.1515 & 0.1664 & 0.1978 & 0.8413 & 0.0839 & 0.5148 \\
 && LLMFactor                & 0.1765 & 0.2008 & 0.2149 & 0.9344 & 0.1039 & 0.5045 \\
 && MarketToM (Qwen2.5-72B)  & 0.2526 & 0.2720 & 0.2488 & 1.0932 & 0.1184 & 0.3805 \\
 && MarketToM (GPT-4o)       & \textbf{0.2625} & \textbf{0.2889} & 0.2484 & \textbf{1.1630} & 0.1169 & 0.3321 \\
\midrule
\multirow{14}{*}{CMIN-US}
 & \multirow{3}{*}{\textbf{ML}}
    & Gaussian NB & -0.0118 & -0.0134 & 0.1788 & -0.0749 & 0.1057 & 0.4655 \\
 && RF           & 0.0395 & 0.0448 & 0.1549 & 0.2892 & 0.0736 & 0.5065 \\
 && Adaboost     & 0.0486 & 0.0532 & 0.1614 & 0.3296 & 0.0688 & 0.5125 \\
\cmidrule(l){2-9}
 & \multirow{3}{*}{\textbf{DL}}
    & RNN  & -0.0096 & -0.0109 & 0.1745 & -0.0625 & 0.0991 & 0.4649 \\
 && LSTM & 0.0511 & 0.0596 & 0.1427 & 0.4177 & 0.0626 & 0.5215 \\
 && GRU  & 0.0226 & 0.0241 & 0.1501 & 0.1606 & 0.0776 & 0.5025 \\
\cmidrule(l){2-9}
 & \multirow{4}{*}{\textbf{MM}}
    & CMIN        & -0.0038 & -0.0043 & 0.1597 & -0.0269 & 0.0962 & 0.5055 \\
 && Melody-GCN  & 0.0485 & 0.0559 & 0.2051 & 0.2725 & 0.0998 & 0.3741 \\
 && ACM24       & 0.0582 & 0.0646 & 0.2083 & 0.3101 & 0.1054 & 0.3485 \\
 && CausalStock & 0.0759 & 0.0825 & 0.1871 & 0.4409 & 0.0842 & 0.4565 \\
\cmidrule(l){2-9}
 & \multirow{4}{*}{\textbf{LLM}}
    & LLM-only                 & 0.0334 & 0.0395 & 0.1815 & 0.2176 & 0.0901 & 0.4381 \\
 && LLMFactor                & 0.0495 & 0.0549 & \textbf{0.1051} & 0.5224 & \textbf{0.0382} & 0.4635 \\
 && MarketToM (Qwen2.5-72B)  & 0.0771 & 0.0832 & 0.1534 & 0.5424 & 0.0797 & 0.3741 \\
 && MarketToM (GPT-4o)       & \textbf{0.1191} & \textbf{0.1288} & 0.2121 & \textbf{0.6073} & 0.1026 & 0.3266 \\
\midrule
\multirow{14}{*}{CMIN-CN}
 & \multirow{3}{*}{\textbf{ML}}
    & Gaussian NB & -0.0446 & -0.0522 & \textbf{0.0751} & -0.6951 & 0.0677 & 0.5025 \\
 && RF           & -0.0235 & -0.0268 & 0.0931 & -0.2879 & 0.0588 & 0.4969 \\
 && Adaboost     & -0.0329 & -0.0383 & 0.0847 & -0.4522 & 0.0574 & 0.4989 \\
\cmidrule(l){2-9}
 & \multirow{3}{*}{\textbf{DL}}
    & RNN  & -0.0351 & -0.0406 & 0.0941 & -0.4315 & 0.0709 & 0.4759 \\
 && LSTM & -0.0312 & -0.0353 & 0.0981 & -0.3598 & 0.0751 & 0.4945 \\
 && GRU  & -0.0132 & -0.0156 & 0.0883 & -0.1767 & 0.0508 & 0.5049 \\
\cmidrule(l){2-9}
 & \multirow{4}{*}{\textbf{MM}}
    & CMIN        &  0.0014 &  0.0016 & 0.0921 &  0.0174 & 0.0503 & 0.4841 \\
 && Melody-GCN  & -0.0131 & -0.0148 & 0.0911 & -0.1625 & 0.0559 & 0.4975 \\
 && ACM24       &  0.0039 &  0.0045 & 0.0975 &  0.0462 & 0.0561 & 0.4925 \\
 && CausalStock &  0.0154 &  0.0176 & 0.0935 &  0.1882 & \textbf{0.0483} & 0.4955 \\
\cmidrule(l){2-9}
 & \multirow{4}{*}{\textbf{LLM}}
    & LLM-only                 & -0.0277 & -0.0324 & 0.0962 & -0.3368 & 0.0751 & 0.5179 \\
 && LLMFactor                & -0.0102 & -0.0116 & 0.0967 & -0.1200 & 0.0692 & 0.4969 \\
 && MarketToM (Qwen2.5-72B)  &  0.0161 &  0.0185 & 0.0915 &  0.2022 & 0.0556 & 0.4815 \\
 && MarketToM (GPT-4o)       & \textbf{0.0403} & \textbf{0.0432} & 0.1341 & \textbf{0.3221} & 0.0717 & 0.4205 \\
\bottomrule
\end{tabular}
}
\end{table}

Table~\ref{tab:trading_simulation} shows that MarketToM variants obtain competitive Sharpe ratios after $10$ bps costs, with the strongest MarketToM variant reaching $1.16$ on ACL18, $0.61$ on CMIN-US, and $0.32$ on CMIN-CN. MarketToM also yields competitive cumulative and annualized returns; notably, GPT-4o remains positive on CMIN-CN while many baselines are negative. Although MarketToM does not uniformly minimize absolute drawdown, its drawdowns remain within the range of the compared methods on CMIN-US and CMIN-CN, and its turnover is comparatively low. These results are consistent with the economic relevance of the classification gains under a simple cost-adjusted long--cash simulation, although the analysis remains an economic sanity check rather than a full trading-strategy study with variable spreads, short selling, market impact, or volatility-adjusted position sizing.
}


\subsection{Qualitative Mechanism Illustrations and Explanatory Value of Second-Order ToM}
\label{sec:validating_cognitive_complexity}
Most LLM-based studies of financial cognition focus only on first-order ToM \citep{yuzhe2025twinmarket,wang2024llmfactor,li2024causalstock}. However, in financial markets, especially in anomaly settings, outcomes are often shaped by recursive cross-group beliefs about others' mental states and anticipated actions, so first-order ToM alone is usually insufficient for mechanism-oriented explanation and may also miss decision-relevant strategic structure. The empirical evidence in this paper comes primarily from the large-sample benchmark comparisons, ablation studies, robustness checks, and statistical-significance tests reported in the main text. Appendix~\ref{app:illustrative_cases} provides two qualitative mechanism illustrations whose purpose is narrower: they show how recursive belief modeling can produce a more behaviorally grounded reading in selected anomaly settings where timing and strategic interaction matter. Both first-order and second-order ToM operate within the same CCN scaffold; the critical distinction lies in whether the model explicitly incorporates recursive inter-agent structures.

\paragraph{Paired Human Evaluation.}
We assess the perceived usefulness of MarketToM's reasoning traces through a within-subjects paired comparison. Fifteen finance-background participants (graduate students and professionals with $\geq$2 years of experience) each evaluate $K=10$ randomly sampled prediction instances covering both anomaly and routine settings. For each instance, they receive anonymized first-order and second-order ToM traces as ``Model~A'' and ``Model~B'' with randomized assignment.

Participants rate each trace on five-point Likert scales for \emph{clarity}, \emph{actionability}, and \emph{trust}, and also answer a forced-choice question on which trace they find more useful for explaining the observed outcome. Likert ratings are averaged within participant across the $K=10$ instances and tested with Wilcoxon signed-rank tests on the 15 paired participant-level scores; forced-choice preferences are computed over all evaluated instances and tested against the $0.5$ random-choice baseline using a binomial test.

\begin{table}[htbp]
\centering
\small
\renewcommand{\arraystretch}{0.8}
\caption{Paired human evaluation of perceived usefulness.}
\begin{tabular}{lccc}
\toprule
\textbf{Dimension} & \textbf{first-order (Mean $\pm$ SD)} & \textbf{second-order (Mean $\pm$ SD)} & \textbf{$p$-value} \\
\midrule
Clarity      & 3.42 $\pm$ 0.61 & 4.11 $\pm$ 0.55 & $<0.001$ \\
Actionability & 3.28 $\pm$ 0.67 & 4.06 $\pm$ 0.58 & $<0.001$ \\
Trust         & 3.35 $\pm$ 0.63 & 4.02 $\pm$ 0.60 & $<0.001$ \\
\midrule
\multicolumn{3}{l}{Forced choice: second-order preferred in 78.7\% of instances (118/150)} & $<0.001$ \\
\bottomrule
\end{tabular}
\label{tab:user_study}
\end{table}

As shown in \Cref{tab:user_study}, second-order reasoning traces receive significantly higher ratings across all three dimensions: clarity (+0.69), actionability (+0.78), and trust (+0.67), all with Wilcoxon signed-rank $p<0.001$. In the forced-choice question, participants select second-order reasoning traces in 78.7\% of evaluated instances (118/150; binomial $p<0.001$). The paired evaluation therefore suggests that the recursive belief structure in second-order ToM is perceived by participants as clearer and more useful for decision support. However, given the modest participant sample and the focus on perceived usefulness rather than explanation faithfulness, the results should be interpreted as preliminary user-facing evidence rather than as validation that the explanations faithfully recover the true market mechanism.

\revise{
\paragraph{External Consistency Analysis of Inferred Mental States.}
We further evaluate inferred belief, intention, and emotion through an external consistency analysis. An independent LLM evaluator receives only the same five days of quote and text information and a fixed rubric, without access to MarketToM's states, prediction, or realized label, and assigns coarse reference labels. MarketToM's open text state descriptions are mapped to the same label space under the identical rubric, after which agreement and Cohen's $\kappa$ are computed by instance, agent role, and mental state dimension.

\begin{table}[htbp]
\centering
\small
\caption{External consistency analysis of inferred mental states.}
\label{tab:external_consistency}
\begin{tabular}{lcc}
\toprule
\textbf{Mental-state dimension} & \textbf{Agreement with reference labels} & \textbf{Cohen's $\kappa$} \\
\midrule
Belief & 73.62\% & 0.5817 \\
Intention & 70.24\% & 0.5324 \\
Emotion & 76.83\% & 0.6213 \\
\midrule
Average & 73.56\% & 0.5785 \\
\bottomrule
\end{tabular}
\end{table}

\Cref{tab:external_consistency} shows agreement above chance across all three dimensions, strongest for emotion and weaker but still positive for intention, which is less directly observable from public quote and text information. Overall, MarketToM's open text state descriptions align systematically with independently assigned labels derived from observable market evidence. These results should be read as partial external consistency evidence, with direct behavioral validation left for future work.
}


\subsection{Ablation Study}
\label{sec:ablation_study}
This paper validates MarketToM through ablations and sensitivity analyses covering the CCN, CEP, second-order ToM, generation temperature $T$, dependency propagation, and prompt format. All ablations use the Qwen2.5-72B backbone under the strict forward protocol and are reported in five parts: \Cref{fig:ablation_results_main}, \Cref{fig:ablation_results_second}, \Cref{tab:tom_counterfactual}, \Cref{tab:ccn_dependency_ablation}, and \Cref{tab:prompt_format_sensitivity}.

\textit{Part I: Main Ablation and Temperature Sensitivity (\Cref{fig:ablation_results_main}).} Robustness to generation randomness is evaluated by varying the LLM generation temperature parameter $T$ across $\{0, 0.7, 1.5\}$. Performance peaks empirically at $T=0.7$, achieving MCC scores of $0.1511$, $0.1529$, and $0.1203$ on ACL18, CMIN-US, and CMIN-CN, respectively. Conversely, a deterministic setting ($T=0$) overly constrains deductive reasoning, reducing MCC to $0.1113$, $0.0647$, and $0.0823$, while an excessively high temperature ($T=1.5$) introduces hallucination noise and yields intermediate MCC values of $0.1394$, $0.1187$, and $0.0978$.

\begin{figure}[pos=t]
\centering
\makebox[\textwidth][c]{%
\includegraphics[width=1.0\linewidth]{figs/ablation_results_main.pdf}%
}
\caption{\linespread{0.85}\selectfont Ablation Results for MarketToM: Components and Temperature Sensitivity.}
\label{fig:ablation_results_main}
\end{figure}

\textit{Part II: Progressive Decomposition (\Cref{fig:ablation_results_second}).}
Performance improves monotonically from MarketToM-NoCEP to MarketToM-1st and then to MarketToM-2nd.
CEP adds MCC improvements of $+0.0420$, $+0.0465$, and $+0.0195$ on ACL18, CMIN-US, and CMIN-CN,
while second-order ToM adds $+0.0315$, $+0.0282$, and $+0.0265$.
Total MCC gains reach $+0.0735$, $+0.0747$, and $+0.0460$, with directionally consistent ACC improvements.

\begin{figure}[pos=t]
\centering
\makebox[\textwidth][c]{%
\includegraphics[width=0.8\textwidth]{figs/ablation_results_second.pdf}%
}
\caption{\linespread{0.85}\selectfont Ablation Results for MarketToM: Incremental Gains.}
\label{fig:ablation_results_second}
\end{figure}

\revise{
\textit{Part III: ToM-Specific Counterfactual Ablation (\Cref{tab:tom_counterfactual}).} The progressive decomposition does not isolate whether second-order gains depend on role-state alignment. We therefore test a role-shuffled variant that keeps the inputs, backbone, CEP library, output format, temperature, and aggregation rule fixed, but randomly permutes inferred mental-state triples across Retail, Institutional, and Arbitrageur agents before action prediction.

\begin{table*}[htbp]
\centering
\small
\renewcommand{\arraystretch}{0.85}
\caption{ToM-specific counterfactual ablation using the Qwen2.5-72B backbone. Values are reported as mean $\pm$ standard deviation over five runs; $\Delta$MCC is computed relative to the full second-order MarketToM on the same dataset.}
\label{tab:tom_counterfactual}
\resizebox{\textwidth}{!}{%
\begin{tabular}{lccccccccc}
\toprule
\multirow{2}{*}{\textbf{Variant}} & \multicolumn{3}{c}{\textbf{ACL18}} & \multicolumn{3}{c}{\textbf{CMIN-US}} & \multicolumn{3}{c}{\textbf{CMIN-CN}} \\
\cmidrule(lr){2-4}\cmidrule(lr){5-7}\cmidrule(lr){8-10}
 & \textbf{ACC} $\uparrow$ & \textbf{MCC} $\uparrow$ & \textbf{$\Delta$MCC} & \textbf{ACC} $\uparrow$ & \textbf{MCC} $\uparrow$ & \textbf{$\Delta$MCC} & \textbf{ACC} $\uparrow$ & \textbf{MCC} $\uparrow$ & \textbf{$\Delta$MCC} \\
\midrule
Full second-order MarketToM & 0.5809 $\pm$ 0.0513 & 0.1511 $\pm$ 0.1160 & 0.0000 & 0.5710 $\pm$ 0.0142 & 0.1529 $\pm$ 0.0362 & 0.0000 & 0.5647 $\pm$ 0.0026 & 0.1203 $\pm$ 0.0159 & 0.0000 \\
First-order only & 0.5631 $\pm$ 0.0490 & 0.1196 $\pm$ 0.1083 & -0.0315 & 0.5582 $\pm$ 0.0150 & 0.1247 $\pm$ 0.0345 & -0.0282 & 0.5521 $\pm$ 0.0031 & 0.0938 $\pm$ 0.0172 & -0.0265 \\
Role-shuffled ToM & 0.5552 $\pm$ 0.0505 & 0.1009 $\pm$ 0.1126 & -0.0502 & 0.5487 $\pm$ 0.0158 & 0.1018 $\pm$ 0.0388 & -0.0511 & 0.5449 $\pm$ 0.0035 & 0.0792 $\pm$ 0.0185 & -0.0411 \\
\bottomrule
\end{tabular}
}
\end{table*}

\Cref{tab:tom_counterfactual} shows that the first-order-only variant loses the gains from second-order recursion, while role shuffling causes larger MCC declines of $0.0502$, $0.0511$, and $0.0411$. Because the market inputs and model configuration are unchanged, the result indicates that role--state alignment contributes to prediction quality.

\revise{
\textit{Part IV: Full Test Set Dependency Ablation (\Cref{tab:ccn_dependency_ablation}).} We evaluate whether CCN-specified dependencies are operationally used through full-test-set dependency ablations. The Qwen2.5-72B backbone, input windows, agent roles, CEP retrieval, temperature, dynamic aggregation, and training-learned CEP library are fixed; only parent-state propagation is perturbed.

The four variants remove $b\rightarrow i$, remove $b\rightarrow e$, remove $\{i,e\}\rightarrow a$, or shuffle belief parents across test instances before downstream inference, preserving the marginal belief-text distribution while breaking instance-level alignment.

\begin{table*}[htbp]
\centering
\small
\renewcommand{\arraystretch}{0.85}
\caption{Full test set dependency ablation of the CCN using the Qwen2.5-72B backbone. Values are reported as mean $\pm$ standard deviation over five runs; $\Delta$MCC is computed relative to the full CCN on the same dataset.}
\label{tab:ccn_dependency_ablation}
\resizebox{\textwidth}{!}{%
\begin{tabular}{lccccccccc}
\toprule
\multirow{2}{*}{\textbf{Variant}} & \multicolumn{3}{c}{\textbf{ACL18}} & \multicolumn{3}{c}{\textbf{CMIN-US}} & \multicolumn{3}{c}{\textbf{CMIN-CN}} \\
\cmidrule(lr){2-4}\cmidrule(lr){5-7}\cmidrule(lr){8-10}
 & \textbf{ACC} $\uparrow$ & \textbf{MCC} $\uparrow$ & \textbf{$\Delta$MCC} & \textbf{ACC} $\uparrow$ & \textbf{MCC} $\uparrow$ & \textbf{$\Delta$MCC} & \textbf{ACC} $\uparrow$ & \textbf{MCC} $\uparrow$ & \textbf{$\Delta$MCC} \\
\midrule
Full CCN & 0.5809 $\pm$ 0.0513 & 0.1511 $\pm$ 0.1160 & 0.0000 & 0.5710 $\pm$ 0.0142 & 0.1529 $\pm$ 0.0362 & 0.0000 & 0.5647 $\pm$ 0.0026 & 0.1203 $\pm$ 0.0159 & 0.0000 \\
w/o $b\rightarrow i$ & 0.5628 $\pm$ 0.0498 & 0.1146 $\pm$ 0.1096 & -0.0365 & 0.5524 $\pm$ 0.0151 & 0.1120 $\pm$ 0.0351 & -0.0409 & 0.5521 $\pm$ 0.0030 & 0.0921 $\pm$ 0.0171 & -0.0282 \\
w/o $b\rightarrow e$ & 0.5659 $\pm$ 0.0502 & 0.1215 $\pm$ 0.1110 & -0.0296 & 0.5561 $\pm$ 0.0148 & 0.1208 $\pm$ 0.0349 & -0.0321 & 0.5530 $\pm$ 0.0029 & 0.0946 $\pm$ 0.0168 & -0.0257 \\
w/o $\{i,e\}\rightarrow a$ & 0.5467 $\pm$ 0.0520 & 0.0860 $\pm$ 0.1141 & -0.0651 & 0.5418 $\pm$ 0.0160 & 0.0837 $\pm$ 0.0381 & -0.0692 & 0.5409 $\pm$ 0.0034 & 0.0710 $\pm$ 0.0188 & -0.0493 \\
Shuffled $b$ parents & 0.5374 $\pm$ 0.0531 & 0.0742 $\pm$ 0.1182 & -0.0769 & 0.5332 $\pm$ 0.0166 & 0.0698 $\pm$ 0.0398 & -0.0831 & 0.5346 $\pm$ 0.0037 & 0.0624 $\pm$ 0.0196 & -0.0579 \\
\bottomrule
\end{tabular}
}
\end{table*}

\Cref{tab:ccn_dependency_ablation} shows that all dependency perturbations reduce performance. Removing $b\rightarrow i$ and $b\rightarrow e$ lowers MCC across all datasets, removing $\{i,e\}\rightarrow a$ causes larger declines, and shuffled belief parents produce the largest average degradation. These results indicate that the CCN functions as an implemented dependency constrained inference scaffold, not only a conceptual diagram, without implying econometric causal identification.
}

\textit{Part V: Prompt-Format Sensitivity (\Cref{tab:prompt_format_sensitivity}).} We test whether performance depends on XML notation or on the broader structured inference design by varying only prompt organization and keeping the rest of the pipeline unchanged.

The HTML/Markdown variant preserves the same semantic fields but changes markup style, while the unstructured variant keeps the task information but removes explicit hierarchical boundaries. All variants retain the same JSON output requirement and Pydantic validation step.

\begin{table*}[htbp]
\centering
\small
\renewcommand{\arraystretch}{0.85}
\caption{Prompt-format sensitivity diagnostic using the Qwen2.5-72B backbone. Values are reported as mean $\pm$ standard deviation over five runs; $\Delta$MCC is computed relative to the XML structured prompt on the same dataset.}
\label{tab:prompt_format_sensitivity}
\resizebox{\textwidth}{!}{%
\begin{tabular}{lccccccccc}
\toprule
\multirow{2}{*}{\textbf{Prompt variant}} & \multicolumn{3}{c}{\textbf{ACL18}} & \multicolumn{3}{c}{\textbf{CMIN-US}} & \multicolumn{3}{c}{\textbf{CMIN-CN}} \\
\cmidrule(lr){2-4}\cmidrule(lr){5-7}\cmidrule(lr){8-10}
 & \textbf{ACC} $\uparrow$ & \textbf{MCC} $\uparrow$ & \textbf{$\Delta$MCC} & \textbf{ACC} $\uparrow$ & \textbf{MCC} $\uparrow$ & \textbf{$\Delta$MCC} & \textbf{ACC} $\uparrow$ & \textbf{MCC} $\uparrow$ & \textbf{$\Delta$MCC} \\
\midrule
XML structured prompt & 0.5809 $\pm$ 0.0513 & 0.1511 $\pm$ 0.1160 & 0.0000 & 0.5710 $\pm$ 0.0142 & 0.1529 $\pm$ 0.0362 & 0.0000 & 0.5647 $\pm$ 0.0026 & 0.1203 $\pm$ 0.0159 & 0.0000 \\
HTML/Markdown structured prompt & 0.5772 $\pm$ 0.0508 & 0.1442 $\pm$ 0.1135 & -0.0069 & 0.5678 $\pm$ 0.0145 & 0.1456 $\pm$ 0.0356 & -0.0073 & 0.5615 $\pm$ 0.0028 & 0.1138 $\pm$ 0.0164 & -0.0065 \\
Unstructured natural-language prompt & 0.5686 $\pm$ 0.0497 & 0.1248 $\pm$ 0.1088 & -0.0263 & 0.5586 $\pm$ 0.0151 & 0.1213 $\pm$ 0.0345 & -0.0316 & 0.5532 $\pm$ 0.0031 & 0.0951 $\pm$ 0.0172 & -0.0252 \\
\bottomrule
\end{tabular}
}
\end{table*}

\Cref{tab:prompt_format_sensitivity} shows that replacing XML tags with HTML/Markdown markers causes only small MCC declines ($0.0065$--$0.0073$), whereas removing explicit prompt boundaries produces larger drops ($0.0252$--$0.0316$). Thus, the results are not tied to XML notation, but structured boundaries help preserve role-specific inputs, CCN parent-state fields, and output constraints.
}


\section{Discussion and Conclusion}\label{sec:conclusion}
Designing effective market-facing DSS remains difficult because financial environments are noisy, adaptive, and strategically interactive \citep{du2024explainable, liu2024multimodal}. MarketToM addresses this challenge by treating forecasting as a cognitively structured decision-support problem rather than a pure pattern-recognition task. The framework combines a ToM-guided Causal Cognitive Network, LLM-based reasoning, and the Cognitive Enhancement Plugin to infer latent mental states from heterogeneous market signals and organize these inferences into explicit reasoning paths. In this design, explanation is generated through the same architecture that produces the prediction, rather than being attached post hoc. Across three financial datasets, MarketToM achieves the strongest overall ACC and MCC among the compared methods, showing that cognitively grounded design can improve both explanatory depth and predictive effectiveness in this application setting.

This paper contributes to the literature in several ways. First, while prior research on AI-enabled financial DSS has largely improved prediction through better signal extraction, richer text representations, or more interpretable variable relations \citep{du2023incorporating, li2024causalstock}, This paper contributes design knowledge by introducing a \textit{Heterogeneous Multi-Agent Framework} with an \textit{Inter-Agent Strategy Refinement} mechanism that supports \textit{second-order} ToM. This shifts the design emphasis from modeling isolated signals to modeling how heterogeneous market participants recursively interpret one another, thereby improving both decision quality and the explanation of financial anomalies. Second, while existing work in predictive analytics and financial text analysis mainly relies on surface text cues, sentiment proxies, or post hoc interpretability \citep{Tetlock2007, fernandez2022explaining}, MarketToM moves toward mechanism-oriented explanation by representing latent mental states in a structured cognitive architecture. In doing so, it helps explain why apparently positive information may already be priced in and why strategic profit-taking can emerge despite favorable news \citep{peterson2002rumor, schmidt2020rumor}. 

Our work offers several practical implications. For asset managers and analysts, MarketToM can function as a scenario-analysis layer that helps distinguish between information that is merely positive in tone and information that is likely to trigger a specific market reaction because it has already been anticipated, doubted, or strategically traded upon. For investor-relations and corporate communication teams, the framework suggests a way to assess how disclosures may be interpreted across audiences, including when apparently favorable messages may still invite skepticism or profit-taking rather than continued buying. More broadly, the study suggests that the practical value of financial DSS may come less from maximizing raw prediction alone and more from helping users reason about when signals are likely to propagate into action, when they may already be priced in, and when the underlying market logic appears fragile or unstable.

Several limitations define the next research agenda. First, the mental state variables in MarketToM are inferred and not directly observed; although the external consistency analysis provides additional evidence, richer behavioral and experimental validation remains necessary. Second, although the framework supports mechanism-oriented interpretation, the current results do not establish strong causal identification. Third, external validity is also constrained by the controlled evaluation setting, including the stock subset and data conditions examined here. \revise{Because the stock universe is filtered by textual coverage, the findings should be interpreted primarily for information-rich stocks, not for all listed firms. In addition, although a cost-adjusted long--cash trading simulation is reported in Section~\ref{sec:trading_simulation}, the current economic evaluation still relies on a single strategy mapping and a fixed $10$ bps transaction-cost assumption, without modeling variable spreads, market impact, or volatility-adjusted position sizing.} Future research should therefore test the framework on broader benchmarks and more heterogeneous market regimes, assess calibration and drift across backbone models, \revise{conduct more systematic sensitivity analyses over alternative look-back windows, trading-oriented target definitions, richer execution-cost models, and pre-specified event-study evaluations,} and extend the design to intraday, multilingual, and cross-asset settings.


\newpage

% \printcredits

%% Loading bibliography style file
% \bibliographystyle{siam}

\bibliographystyle{cas-model2-names}
\let\oldbibliography\thebibliography
\renewcommand{\thebibliography}[1]{%
\oldbibliography{#1}%
\baselineskip8pt 
\setlength{\itemsep}{3pt}
}
\bibliography{bib}

% Loading bibliography database

% Biography
%\bio{}
% Here goes the biography details.
%\endbio

%\bio{pic1}
% Here goes the biography details.
%\endbio
% \bibliography{bib}


\newpage
\appendix


\section{Benchmark Implementation Details}
\label{app:benchmark_implementation}
This appendix documents the comparability-relevant benchmark protocol and does not attempt to list every software option. All methods use the same underlying five-day quote--text information source, next-day label definition, repeated five-stock draws, train/validation/test splits, and held-out instances. Hyperparameters are selected on validation data using common search spaces across datasets; stochastic and LLM-dependent methods are averaged over five runs.

Table~\ref{tab:app_benchmark_settings} summarizes the aligned inputs, LLM usage, and representative settings for each baseline family. \revise{Full prompt templates and the runnable configuration are provided in the project repository.\footnote{\url{https://github.com/yyveggie/MarketToM}}}

\begin{table}[b]
\centering
\scriptsize
\renewcommand{\arraystretch}{0.95}
\caption{Representative benchmark implementation settings. The table reports the main comparability-relevant settings, not every implementation option.}
\begin{tabularx}{\textwidth}{>{\raggedright\arraybackslash}p{2.0cm}>{\raggedright\arraybackslash}p{2.55cm}>{\raggedright\arraybackslash}p{2.55cm}>{\raggedright\arraybackslash}X}
\toprule
\textbf{Method} & \textbf{Aligned inputs} & \textbf{LLM usage} & \textbf{Representative settings} \\
\midrule
Gaussian NB / RF / AdaBoost & Five-day quote features and text embeddings projected into a common latent representation. & None. & Gaussian NB uses no grid search. RF uses 500 trees, max depth 10, \texttt{sqrt} feature sampling, and minimum leaf size 5. AdaBoost is tuned over weak learners, learning rate, and tree depth; the selected setting uses 200 weak learners, learning rate 0.1, and depth-2 decision trees. \\
\midrule
RNN / LSTM / GRU & Same five-day sequence with concatenated quote/text latent input. & None. & Shared validation grid over hidden size $\{32,64,128\}$, recurrent layers $\{1,2\}$, dropout $\{0,0.2,0.5\}$, and learning rate $\{10^{-4},5{\times}10^{-4},10^{-3}\}$. Selected settings use one layer, hidden size 64, dropout 0.2, Adam, batch size 64, early stopping patience 10, and cross-entropy loss; learning rate is $10^{-3}$ for RNN/GRU and $5{\times}10^{-4}$ for LSTM. \\
\midrule
CMIN & Text and high/low/close price inputs, with cross-stock causal matrix estimated from prices. & None. & Uses five-day history, GloVe/Bi-GRU-style text encoder in the original architecture, memory depth 3, price/text hidden size 50, dropout 0.2, Adam, and binary cross-entropy/classification training. The causal matrix is computed from price sequences rather than from an external industry graph or knowledge graph. \\
\midrule
ACM24 & Text and OHLC/adjusted-close-derived quantitative indicators, with batch-internal contrastive sampling. & None; Sentence-BERT is used as an encoder-only text representation model in the original design. & Uses the aligned five-day window, 11 quantitative indicators, Adam with learning rate $10^{-6}$, 100 epochs, hidden dimension 256, triplet top-$k$ selection from $\{10,20\}$, and triplet margin 1.0 when the original paper does not specify the margin. \\
\midrule
Melody-GCN & Quote/text inputs plus method-specific technical indicators and learned dynamic relation matrices. & None; BERT-style encoder features are used for text representation, not generation. & Uses 11 price-side inputs, 768-dimensional text features, dropout 0.1, batch size 16, Adam with learning rate $10^{-4}$, and 500 epochs. For the unified five-day protocol, temporal scales and internal adjacency dimensions are adapted to the five-stock subset while preserving the dynamic graph-convolution design. \\
\midrule
CausalStock & Price and news-text signals; no external fixed relation graph in the main setting. & Yes, only for offline denoised-news scoring; standardized to GPT-4o in our benchmark. & Uses max lag 5, batch size 32, learning rate $10^{-5}$, price hidden dimension 4, news embedding dimension 64, sparse-graph penalty 1.0, BCE weight 0.01, and a 3-layer FCM. The LLM converts each news item into five scores, which are then embedded by a linear layer; it does not directly perform final prediction. \\
\midrule
LLM-only / LLMFactor & Same five-day quote--text window supplied in prompt form. & GPT-4o, temperature $0$, no external retrieval. & LLM-only is zero-shot and outputs only \textit{Buy} or \textit{Sell}; its prompt excludes MarketToM-specific ToM and agentic structures. LLMFactor follows the factor-extraction prompting design while using the same aligned input window and backend. \\
\bottomrule
\end{tabularx}
\label{tab:app_benchmark_settings}
\end{table}


\section{Illustrative Mechanism Cases}
\label{app:illustrative_cases}
This appendix provides qualitative mechanism examples for interpreting second-order ToM in selected anomaly settings. These cases are intended to complement, rather than replace, the large-sample benchmark comparisons, ablation studies, robustness checks, and statistical-significance tests reported in the main text. Their role is to make the model's reasoning process transparent in representative anomaly settings and to illustrate how second-order ToM can support mechanism-oriented interpretation.


\subsection{Illustrative Anomaly 1: ``Buy the Rumor, Sell the News'' (BTRSTN)}
We select the BTRSTN illustration using transparent heuristic rules: a pre-run-up, a hollow rally before the focal news (return $\geq +4\%$ and $VR=\text{Vol}/MA20 \leq 0.95$), and an event-day reversal (return $\leq -3\%$) despite non-negative event sentiment. The same $[t-1,t+1]$ news screen excludes windows dominated by firm-specific negative catalysts or broad market shocks; the selected Tesla window in Figure~\ref{fig:btrstn_case} shows a $+5.45\%$ hollow rally on October~6,~2025 followed by a $-4.45\%$ decline on October~7, so it is used as an illustrative mechanism example rather than a large-sample event-study identifier.

\begin{figure}[pos=t]
\centering
\includegraphics[width=0.9\columnwidth]{figs/case_study_tesla_dual.pdf}
\caption{Illustrative BTRSTN mechanism example for Tesla (Sep 16--Oct 10, 2025).}
\label{fig:btrstn_case}
\end{figure}

As illustrated in \Cref{fig:tesla_case_path}, under first-order ToM, the Institutional agent infers Retail's market-facing mental states: Retail interprets the Model~Y/3 unveiling as bullish and tends to accumulate. Conditioned on the observed environmental states and the retrieved first-order strategies, this yields a plausible first-order reading of the selected window along the CCN pathway. However, this account remains behaviorally incomplete: it does not explicitly represent recursive inter-agent beliefs, and therefore does not by itself explain why Retail demand may persist strongly enough to serve as an exit channel, nor why Oct~7 is a plausible distribution window.

\begin{figure}[pos=t]
\centering
\includegraphics[width=1.0\linewidth]{figs/tesla_case.pdf}
\caption{Summarized mechanism trace for the Tesla BTRSTN case.}
\label{fig:tesla_case_path}
\end{figure}

Second-order ToM adds the missing strategic layer. Beyond inferring Retail's bullish market belief, the Institutional agent infers \emph{Retail's belief about Institutional behavior}: Retail observes the Oct~6 price rally and interprets it as evidence that Institutions are supporting the move and will continue to do so. The price--volume divergence then suggests that this belief may be a misperception, because the Oct~6 rally occurs on relatively thin volume, more consistent with limited Institutional participation than broad accumulation. Under this interpretation, Retail may continue buying into the product-news window precisely because it believes Institutions are behind the rally; this recursive-belief mismatch creates a plausible liquidity window into which Institutions may distribute on Oct~7.

The first-order ToM variant (MarketToM-1st) captures the first-order layer, Retail's bullish market belief, but lacks the second-order component needed to account for \emph{why} Retail demand may remain an exit channel at this specific juncture. In this illustrative mechanism case, both MarketToM-1st and MarketToM-2nd (the full MarketToM) produce a Sell prediction, but the second-order ToM reasoning yields a more decision-informative account of \emph{why} Sell is warranted here: it links the action to a fragile recursive-belief liquidity window, a distribution-oriented intention, and a disciplined vigilant emotional state. Accordingly, the second-order structure does not change the directional prediction in this case; rather, it sharpens the mechanism account of why Sell is warranted at this specific juncture, aligning the final action more closely with behavioral-finance discussions of BTRSTN~\citep{daniel1998,deBondt1985}.


\subsection{Illustrative Anomaly 2: Riding the Bubble via ToM}
A fundamental puzzle in behavioral finance is why investors frequently choose to buy assets at prices they know exceed fundamental values, a phenomenon known as ``riding the bubble.'' Neuroeconomic evidence suggests that participants with stronger mentalizing abilities are often more deeply involved in bubble dynamics \citep{bossaerts2019perception,de2013mind,hefti2025mental}. However, the ToM evaluated in these studies primarily reflects general mentalizing tendency and intensity, rather than explicitly distinguishing between first-order and second-order mental recursion.

An illustrative vignette is constructed to show how MarketToM's second-order ToM design may refine this distinction. It is set during a bubble period and zooms in on a single-day snapshot near the top. Following \citet{de2013mind}, the asset's market price ($\$100$) is set to deviate sharply from its fundamental value ($\$20$), accompanied by textual sentiment indicating strong retail buying enthusiasm. To isolate the ToM effect, this static vignette bypasses the CEP and relies entirely on direct mental evolution. Figure~\ref{fig:bubble_case_path} summarizes the resulting mechanism trace.

Under first-order ToM, the institutional agent infers strong bullishness among retail investors but does not explicitly model how that demand is recursively sustained or when it may break. As a result, the agent may overweight visible buying pressure relative to valuation in the immediate decision, forming a trend-following intention and a speculative emotion that together support a Buy action. In this illustrative setting, first-order reasoning acts less as a pure fundamental arbitrage logic than as a contributor to anomaly continuation.

Second-order ToM operates on the same CCN pathway but adds the missing strategic layer: besides inferring retail's bullishness, the institutional agent infers \emph{retail's recursive belief about other market participants}, namely that continued buying is sustained by expectations of further liquidity from others. Near the top of the bubble, this recursive belief loop is more informative as fragile exit liquidity than as evidence of durable upside. Accordingly, the agent shifts to a profit-taking intention and a vigilant, opportunistic emotional state, causing the final action to flip to Sell. The advantage of second-order ToM reasoning in this vignette is therefore timing: it recognizes the loop as exploitable but fragile and exits one step earlier.

Taken together, the contrast suggests that first-order ToM can be read as following anomaly continuation, whereas second-order ToM may better capture the timing logic needed to exit fragile recursive-belief loops. First-order reasoning follows visible momentum; second-order reasoning recognizes that the same momentum may be sustained by a fragile recursive-belief structure and interprets that structure as a timing risk. This vignette therefore illustrates how second-order ToM can provide a mechanism-oriented account of timing decisions near a presumed bubble peak, complementing the large-sample predictive and ablation evidence reported in the main analysis.

\begin{figure}[pos=t]
\centering
\includegraphics[width=1.0\columnwidth]{figs/bubble.pdf}
\caption{Summarized mechanism trace for the ``Riding the Bubble'' vignette.}
\label{fig:bubble_case_path}
\end{figure}


\section{Additional Diebold--Mariano Pairwise Tests}
\label{app:dm_tests}
This appendix reports representative paired Diebold--Mariano (DM) tests for MarketToM (GPT-4o) against strong baselines under 0--1 loss. Predictions are matched at the sample level, and the loss difference is defined as baseline loss minus MarketToM loss, so a positive value indicates lower MarketToM loss.

\begin{table}[htbp]
\centering
\small
\setlength{\tabcolsep}{4pt}
\renewcommand{\arraystretch}{0.88}
\caption{Representative pairwise Diebold--Mariano tests for MarketToM (GPT-4o).}
\begin{tabular}{llrrr}
\toprule
\textbf{Dataset} & \textbf{Baseline} & \textbf{Loss Diff.} & \textbf{DM Stat.} & \textbf{Two-sided $p$} \\
\midrule
\multirow{6}{*}{ACL18}
& ACM24 & 0.0660 & 2.340 & 0.019 \\
& LLM-only & 0.0615 & 2.279 & 0.023 \\
& Melody-GCN & 0.0540 & 1.843 & 0.065 \\
& LLMFactor & 0.0455 & 1.612 & 0.107 \\
& CausalStock & 0.0325 & 1.114 & 0.265 \\
& CMIN & 0.0250 & 0.894 & 0.371 \\
\midrule
\multirow{6}{*}{CMIN-US}
& CMIN & 0.0710 & 4.080 & $<0.001$ \\
& LLM-only & 0.0625 & 3.158 & 0.002 \\
& Melody-GCN & 0.0415 & 2.197 & 0.028 \\
& ACM24 & 0.0370 & 2.062 & 0.039 \\
& LLMFactor & 0.0285 & 1.673 & 0.094 \\
& CausalStock & 0.0305 & 1.598 & 0.110 \\
\midrule
\multirow{6}{*}{CMIN-CN}
& ACM24 & 0.0605 & 3.269 & 0.001 \\
& CMIN & 0.0585 & 3.165 & 0.002 \\
& LLM-only & 0.0565 & 3.152 & 0.002 \\
& Melody-GCN & 0.0485 & 2.868 & 0.004 \\
& LLMFactor & 0.0475 & 2.663 & 0.008 \\
& CausalStock & 0.0310 & 1.744 & 0.081 \\
\bottomrule
\end{tabular}
\begin{flushleft}
\footnotesize Note: Comparisons are paired on the same held-out test predictions under 0--1 loss. Positive loss differences indicate lower MarketToM loss. The attained significance reflects the pooled held-out sample size of each dataset, so some comparisons on the smaller ACL18 test set are positive but not significant under the conservative two-sided test.
\end{flushleft}
\label{tab:app_dm_pairwise_strong_baselines}
\end{table}


\end{document}